% Options for packages loaded elsewhere
\PassOptionsToPackage{unicode}{hyperref}
\PassOptionsToPackage{hyphens}{url}
\documentclass[
  ukrainian,
]{article}
\usepackage{xcolor}
\usepackage[margin=2.5cm]{geometry}
\usepackage{amsmath,amssymb}
\setcounter{secnumdepth}{-\maxdimen} % remove section numbering
\usepackage{iftex}
\ifPDFTeX
  \usepackage[T1]{fontenc}
  \usepackage[utf8]{inputenc}
  \usepackage{textcomp} % provide euro and other symbols
\else % if luatex or xetex
  \usepackage{unicode-math} % this also loads fontspec
  \defaultfontfeatures{Scale=MatchLowercase}
  \defaultfontfeatures[\rmfamily]{Ligatures=TeX,Scale=1}
\fi
\usepackage{lmodern}
\ifPDFTeX\else
\setmainfont{DejaVu Serif}
\setsansfont{DejaVu Sans}
\setmonofont{DejaVu Sans Mono}
\setmathfont{STIX Two Math}
\fi
\ifPDFTeX\else
  % xetex/luatex font selection
\fi
% Use upquote if available, for straight quotes in verbatim environments
\IfFileExists{upquote.sty}{\usepackage{upquote}}{}
\IfFileExists{microtype.sty}{% use microtype if available
  \usepackage[]{microtype}
  \UseMicrotypeSet[protrusion]{basicmath} % disable protrusion for tt fonts
}{}
\makeatletter
\@ifundefined{KOMAClassName}{% if non-KOMA class
  \IfFileExists{parskip.sty}{%
    \usepackage{parskip}
  }{% else
    \setlength{\parindent}{0pt}
    \setlength{\parskip}{6pt plus 2pt minus 1pt}}
}{% if KOMA class
  \KOMAoptions{parskip=half}}
\makeatother
\usepackage{longtable,booktabs,array}
\newcounter{none} % for unnumbered tables
\usepackage{calc} % for calculating minipage widths
% Correct order of tables after \paragraph or \subparagraph
\usepackage{etoolbox}
\makeatletter
\patchcmd\longtable{\par}{\if@noskipsec\mbox{}\fi\par}{}{}
\makeatother
% Allow footnotes in longtable head/foot
\IfFileExists{footnotehyper.sty}{\usepackage{footnotehyper}}{\usepackage{footnote}}
\makesavenoteenv{longtable}
\usepackage{graphicx}
\graphicspath{{/Users/kalinicenkovitalijmikolajovic/.codex/worktrees/5c15/Script_for_translate_book/Output_Final/}}
\setkeys{Gin}{width=0.82\linewidth,height=0.82\textheight,keepaspectratio}
\makeatletter
\newsavebox\pandoc@box
\newcommand*\pandocbounded[1]{% scales image to fit in text height/width
  \sbox\pandoc@box{#1}%
  \Gscale@div\@tempa{\textheight}{\dimexpr\ht\pandoc@box+\dp\pandoc@box\relax}%
  \Gscale@div\@tempb{\linewidth}{\wd\pandoc@box}%
  \ifdim\@tempb\p@<\@tempa\p@\let\@tempa\@tempb\fi% select the smaller of both
  \ifdim\@tempa\p@<\p@\scalebox{\@tempa}{\usebox\pandoc@box}%
  \else\usebox{\pandoc@box}%
  \fi%
}
% Set default figure placement to htbp
\def\fps@figure{htbp}
\makeatother
\ifLuaTeX
\usepackage[bidi=basic,shorthands=off,provide=*]{babel}
\else
\usepackage[bidi=default,shorthands=off,provide=*]{babel}
\fi
\ifLuaTeX
  \usepackage{selnolig} % disable illegal ligatures
\fi
\setlength{\emergencystretch}{3em} % prevent overfull lines
\providecommand{\tightlist}{%
  \setlength{\itemsep}{0pt}\setlength{\parskip}{0pt}}
\usepackage{bookmark}
\IfFileExists{xurl.sty}{\usepackage{xurl}}{} % add URL line breaks if available
\urlstyle{same}
\hypersetup{
  pdflang={uk},
  hidelinks,
  pdfcreator={LaTeX via pandoc}}

\author{}
\date{}

\begin{document}

\pandocbounded{\includegraphics[keepaspectratio]{images/_page_353_Figure_9.jpeg}}

\texttt{У\ конформному\ відображенні\ рівняння\ Лапласа\ \textquotesingle{}\textbackslash{}varphi\_\{xx\}+\textbackslash{}varphi\_\{yy\}=0\textquotesingle{}\ у\ координатній\ площині\ (x,y)\ повертається\ у\ рівняння\ Лапласа\ \textquotesingle{}\textbackslash{}varphi\_\{nn\}\textbackslash{}dashv\ -\textbackslash{}varphi\_\{vo\}=0\textquotesingle{}\ у\ координатній\ площині\ (u,v).\ Інші\ рівняння,\ такі\ як\ \textquotesingle{}\textbackslash{}varphi\_\{xx\}+2\textbackslash{}varphi\_\{xy\}+\textbackslash{}varphi\_\{yy\}=0\textquotesingle{}\ ,\ не\ мають\ такої\ властивості.\ Після\ розв\textquotesingle{}язання\ рівняння\ Лапласа\ \textquotesingle{}\textbackslash{}varphi(u,v)\textquotesingle{}\ у\ простій\ області\ (наприклад,\ у\ колі,\ напівплощині,\ квадраті)\ знайдено,\ достатньо\ підставити\ вирази\ U=\textquotesingle{}U(X,\ Y)\textquotesingle{},\ V=\textquotesingle{}V(X,\ Y)\textquotesingle{}\ у\ цей\ розв\textquotesingle{}язок,\ і\ ми\ отримаємо\ розв\textquotesingle{}язок\ нашої\ задачі\ \textquotesingle{}\textbackslash{}phi{[}u(x,\ y),\ v(x,\ y){]}\textquotesingle{}\ ,\ вираженого\ через\ початкові\ змінні.У\ цій\ лекції\ ми\ розглянемо:1)\ про\ властивості\ конформних\ відображень,2)\ як\ їх\ будувати\ (для\ деяких\ районів),3)\ як\ розв\textquotesingle{}язувати\ задачі\ методом\ конформних\ відображень.Однією\ з\ головних\ труднощів,\ що\ виникають\ при\ розв\textquotesingle{}язанні\ задач\ меж,\ є\ складна\ форма\ кордонів.\ Навіть\ межі\ відносно\ простої\ форми\ часто\ роблять\ задачу\ дуже\ складною\ для\ розв\textquotesingle{}язання.\ Звісно,\ можна\ спробувати\ розв\textquotesingle{}язати\ задачу\ за\ допомогою\ методів\ теорії\ збурень,\ але\ теорія\ збурень\ працює\ добре\ лише\ у\ випадках,\ коли\ межа\ близька\ до\ простої.Однак\ існує\ спосіб\ розв\textquotesingle{}язати\ двовимірне\ рівняння\ Лапласа\ в\ областях\ з\ комплексною\ межею.\ Цей\ метод\ базується\ на\ конформних\ відображеннях.Перш\ ніж\ почати\ розв\textquotesingle{}язувати\ задачі\ цим\ методом,\ давайте\ трохи\ ознайомимося\ з\ поняттям\ конформного\ відображення\ та\ комплексних\ змінних\ загалом.\#\ Комплексні\ функції\ та\ конформні\ відображенняУ\ цій\ лекції\ ми\ ознайомимося\ лише\ з\ тими\ поняттями\ з\ теорії\ функцій\ комплексної\ змінної,\ які\ нам\ знадобляться\ пізніше.\ Представимо\ комплексне\ число\ z\ =\ x\ +\ iy\ як\ точку\ на\ площині\ комплексного\ xy\ (рис.\ 47.1).Щоб\ ввести\ поняття\ конформного\ відображення,\ спершу\ ознайомимося\ з\ поняттям\ функції\ комплексної\ змінної}

\begin{verbatim}
.Тут z — це комплексна змінна (з деякої області площини z), а w — нова комплексна змінна, яка отримає значення відповідно до формули w = f(z) у новій комплексній площині w.Наприклад, комплексна функція, визначена в *першому квадранті* комплексної z-площини, w=z^2 відображає цю область у всю півплощину v>0 у комплексній площині w.Таблиця 47 2| Площина d | z \longrightarrow z^2 | Площина S |
\end{verbatim}

\textbar--------------------------\textbar-------------------------\textbar------------------------------------------------------------\textbar{}

0 \textbar{} \textbar{} 0 \textbar{}

t1 → t \textbar{} \textbar{} -1(2) \textbar{}

Вісь x Перший квадрант \textbar{} \textbar{} Позитивна дійсна піввісь
верхньої півплощини \textbar Таблиця 47.2 показує відповідність, яка
встановлюється шляхом відображення w=z\^{}2 між деякими точками на
площинах z і w.

\pandocbounded{\includegraphics[keepaspectratio]{images/_page_355_Figure_5.jpeg}}

\begin{verbatim}
Рис. 47.1. Площина комплексної змінної \boldsymbol{z} та некольорових формул бігунів Р И С У Н О К
\end{verbatim}

\pandocbounded{\includegraphics[keepaspectratio]{images/_page_355_Figure_7.jpeg}}

\begin{verbatim}
47.2. Відображення першого квадранта площини {\bf z} у верхню півплощину {\bf w} .Інакше кажучи, якщо ми говоримо про комплексісні функції змінної w = f(z), то можна розглянути питання, як криві площини z перетворюються на криві в площиніЦе питання не ставиться в теорії функцій дійсної змінної, з якою читач уже добре знайомий.Перш ніж рухатися далі, давайте розглянемо, як w=z^2 складне відображення типів у еквівалентній дійсній формі. Дійсна форма показує, як початкові координати (x,y) у площині z перетворюються на нові координати (u,v) у площині w w=z^2 .
\end{verbatim}

\begin{verbatim}
\begin{cases} u(x, 0) = x, \\ u_t(x, 0) = 0. \end{cases}
\end{verbatim}

І давайте окремо прирівняємо реальні та уявні частини в цьому
співвідношенні. У результаті ми отримуємо

\begin{verbatim}
\begin{array}{ll} u=x^2-y^2,\\ v=2xy, \end{array}
\end{verbatim}

\begin{verbatim}
(действительная форма отображения w=z^2 ).Ця форма нам знадобиться пізніше. У цьому прикладі це цікаво, оскільки показує, як гіперболи в площині z переходять до координатних прямих u = const і v = const у площині w.Тепер перейдемо до вивчення одного конкретного типу функцій комплексної змінної, які виконують конформні відображення.# Определение конформного отображенияВідображення w=f(z) комплексної площини z у комплексну площину w називається конформним у точці z_{\bullet} площини z, якщо похідна f'(z_{\bullet}) \neq 0 . Відображення f(z) називається конформним у регіоні D, якщо f'(z) \neq 0 в кожній точці регіону D.Наприклад, відображення f(z) = z^2 конформно всюди, крім z = 0, оскільки z'(z) = 2z \neq 0 для всіх z \neq 0 . З іншого боку, відображення e^z конформно по всій площині z, оскільки скрізь f'(z) = e^z \neq 0 . Яке призначення конформних відображень? Відповідь полягає в тому, що при розв'язанні рівняння Лапласа \phi_{xx} + \phi_{yy} = 0 в межах деякої області площини змінних x і y, ми можемо розглядати цю площину як площину комплексної змінної z. Розглянемо комплексне відображення w = f(z) площини z у площину w, у якій введені координати (u, v). Коли w = f(z \phi_{xy} + \phi_{uy} = 0 ) відображається, рівняння Лапласа \varphi_{xx} + \varphi_{yy} = 0 перетворюється у нове диференціальне рівняння в похідних похідних, яке залежить від нових координат і v. To — рівняння, яке \phi_{xx} + \phi_{uu} = 0 повернеться у рівняння \varphi_{nn} + \varphi_{nn} = 0 . Таким чином, ми придумаємо ідею знайти конформне відображення, яке переносить область з комплексною межею у область з простою межею (нагадаємо, що відображення w=z^* переносить межу першого квадранта площини z у фактичну вісь площини w).Ось кілька прикладів, які покажуть, як можна розв'язати диференціальні рівняння у похідних за допомогою конформних відображень.Уравнение Лапласа в верхней полуплоскости.Припустимо, що нам потрібно розв'язати наступну задачу Діріхле у верхній половині видимості (рис. 47.3):(47.1)
\end{verbatim}

\texttt{можно\ найти\ стандартным\ методом\ разделения\ переменных.\ При\ этом\ оказывается,\ что\ решение\ \ \textquotesingle{}u\textbackslash{}left(x,\textbackslash{},t\textbackslash{}right)\textquotesingle{}\ ,\ представленное\ в\ виде\ ряда}

Один із способів розв'язання цієї проблеми --- застосувати перетворення
Фур'є до змінної x, але набагато краще конформно відобразити верхню
половину площини у площину (u, v). Читачеві не буде складно побачити, що
конформне відображення

является суперпозицией простейших колебаний. Пространственная часть
\texttt{X\_n(x)} каждого элементарного колебания является решением
(собственной функцией) некоторой краевой задачи Штурма---Лиувилля.Итак,
мы научились решать волиовое уравнение
\texttt{u\_\{II\}=c\^{}2u\_\{xx\}} на всей прямой
\texttt{-\textbackslash{}infty\ \textless{}\ x\ \textless{}\ \textbackslash{}infty}
. Формула Даламбера представляет решение в виде суммы двух бегущих волн,
распространяющихся в противоположные стороны. Если же рассматривать это
же уравнение на ограниченном промежутке 0 \textless{} x \textless{} L,
то бегущих волн уже не станет, так как они будут взаимодействовать с
границами. Вместо них возникают другие волны, которые называют стоячими
волнами. Посмотрим, например, что случится, если гитарную струну,
закрепленную в точках x=0, x=L, привести в движение. Для ответа на этот
вопрос нам необходимо решить следующую смешанную задачу гиперболического
типа:

\begin{verbatim}
Вона має такі властивості (див. RNS. 47.4).1. Вона відображає верхню півплощину площини z у смуги -\infty < u < \infty , 0 < v < \pi , у площині w.
\end{verbatim}

\pandocbounded{\includegraphics[keepaspectratio]{images/_page_357_Figure_11.jpeg}}

\begin{verbatim}
Фіг. 47.3 Проблема граничного значення у верхній половині площі.- 2. Ребро прямої y=0, -1 < x < 1, у площині z (де потенціал дорівнює 1), переходить у пряму v=\pi, - \infty < u < \infty , у площині w. - 3. Два промені x>1, y=0, і x<-1, y=0, у площині z проходять відповідно у позитивний і негативнийЧитач повинен бути впевнений, що може підтвердити ці твердження.
\end{verbatim}

\pandocbounded{\includegraphics[keepaspectratio]{images/_page_358_Figure_2.jpeg}}

Рис. 47.4. Перетворення складної задачі на легку за допомогою
конформного відображення.Важливість цих властивостей полягає в тому, що
вони перетворюють початкову вадаху (47.1) на дуже просту задачу(47.2)

Розв'язок очевидний:
\texttt{\textbackslash{}varphi(u,\ v)\ =\ \textbackslash{}frac\{1\}\{\textbackslash{}pi\}v}
.Тепер, щоб знайти розв'язок початкової задачі, потрібно пояснити v
через x і y і підставити цей вираз на
\texttt{q(u,v)=v\textquotesingle{}\textbackslash{}pi} . Після цього ми
отримуємо

Виявляється, що в цьому випадку хвилі, що рухаються, відбиваються від
меж так, що отримані коливання стають непоширюючими, а зберігають форму
в одному положенні, тобто Ф і г

\pandocbounded{\includegraphics[keepaspectratio]{images/_page_359_Figure_10.jpeg}}

47.5. Конформне відображення на кільце.Це вже не так очевидно, але
існують довідники, де можна знайти дисплей із вказаними властивостями.
Одним із найзагальних є «Довідник конформних відображень»G. Kober 1),
яка містить сотні конформних відображень між різними регіонами. Якщо
читач витратить трохи часу і краще ознайомиться з конформними таблицями
відображення, він знайде відповідне конформне відображення для задачі
(47.3). У нашому конкретному випадку це буде образ виду(47.4)

де s=-0,146, а t=-6,85. У своїй еквівалентній дійсній формі це
буде(47.5)

\begin{verbatim}
\begin{cases} u(x, 0) = f(x), \\ u_t(x, 0) = g(x), \end{cases}
v = \gamma [y(x-t)],
\gamma = 2t/[(x-t)^2 + y^2].
\end{verbatim}

найдем сначала стоячне волны, т. е. решения вида

\begin{verbatim}
\begin{aligned} \varphi_{uu} + \varphi_{vv} &= 0, & \text{в кольче,} \\ (\Gamma Y) & & & & & & & \\ (\Gamma Y) & & & & & & \\ (\Gamma Y) & & & & & \\ & & & & & \\ (\Gamma Y) & & & & & \\ & & & & \\ & & & & \\ & & & & \\ & & & & \\ & & & \\ & & & \\ & & & \\ & & & \\ & & & \\ & & & \\ & & \\ & & \\ & & \\ & & \\ & & \\ & & \\ & & \\ & & \\ & & \\ & & \\ & & \\ & & \\ & & \\ & & \\ & & \\ & & \\ & & \\ & & \\ & & \\ & & \\ & & \\ & & \\ & & \\ & & \\ & & \\ & & \\ & & \\ & & \\ & & \\ & & \\ & & \\ & & \\ & & \\ & & \\ & & \\ & & \\ & & \\ & & \\ & & \\ & & \\ & & \\ & & \\ & & \\ & & \\ & & \\ & & \\ & & \\ & & \\ & & \\ & & \\ & & \\ & & \\ & & \\ & & \\ & & \\ & & \\ & & \\ & & \\ & & \\ & & \\ & & \\ & & \\ & & \\ & & \\ & & \\ & & \\ & & \\ & & \\ & & \\ & & \\ & & \\ & & \\ & & \\ & & \\ & & \\ & & \\ & & \\ & & \\ & & \\ & & \\ & & \\ & & \\ & & \\ & & \\ & & \\ & & \\ & & \\ & & \\ & & \\ & & \\ & & \\ & & \\ & & \\ & & \\ & & \\ & & \\ & & \\ & & \\ & & \\ & & \\ & & \\ & & \\ & & \\ & & \\ & & \\ & & \\ & & \\ & & \\ & & \\ & & \\ & & \\ & & \\ & & \\ & & \\ & & \\ & & \\ & & \\ & & \\ & & \\ & & \\ & & \\ & & \\ & & \\ & & \\ & & \\ & & \\ & & \\ & & \\ & & \\ & & \\ & & \\ & & \\ & & \\ & & \\ & & \\ & & \\ & & \\ & & \\ & & \\ & & \\ & & \\ & & \\ & & \\ & & \\ & & \\ & & \\ & & \\ & & \\ & & \\ & & \\ & & \\ & & \\ & & \\ & & \\ & & \\ & & \\ & & \\ & & \\ & & \\ & & \\ & & \\ & & \\ & & \\ & & \\ & & \\ & & \\ & & \\ & & \\ & & \\ & & \\ & & \\ & & \\ & & \\ & & \\ & & \\ & & \\ & & \\ & & \\ & & \\ & & \\ & & \\ & & \\ & & \\ & & \\ & & \\ & & \\ & & \\ & & \\ & & \\ & & \\ & & \\ & & \\ & & \\ & & \\ & & \\ & & \\ & & \\ & & \\ & & \\ & & \\ & & \\ & & \\ & & \\ & & \\ & & \\ & & \\ & & \\ & & \\ & & \\ & & \\ & & \\ & & \\ & & \\ & & \\ & & \\ & & \\ & & \\ & & \\ & & \\ & & \\ & & \\ & & \\ & & \\ & & \\ & & \\ & & \\ & & \\ & & \\ & & \\ & & \\ & & \\ & & \\ & & \\ & & \\ & & \\ & & \\ & & \\ & & \\ & & \\ & & \\ & & \\ & & \\ & & \\ & & \\ & & \\ & & \\ & & \\ & & \\ & & \\ & & \\ & & \\ & & \\ & & \\ & & \\ & & \\ & & \\ & & \\ & & \\ & & \\ & & \\ & & \\ & & \\ & & \\ & & \\ & & \\ & & \\ & & \\ & & \\ & & \\ & & \\ & & \\ & & \\ & & \\ & & \\ & & \\ & & \\ & & \\ & & \\ & & \\ & & \\ & & \\ & & \\ & & \\ & & \\ & & \\ & & \\ & & \\ & & \\ & & \\ & & \\ & & \\ & & \\ & & \\ & & \\ & & \\ & & \\ & & \\ & & \\ & & \\ & & \\ & & \\ & & \\ & & \end{aligned}
\end{verbatim}

Підставляючи цей вираз у рівняння (20.1) і ділячи змінні, отримуємо дві
звичайні диференціальні рівняння

\begin{verbatim}
\varphi(r) = a \ln r + b, \quad r = u^2 + v^2.
\varphi(u, v) = 0.57 \ln(u^2 + v^2) + 1
\end{verbatim}

\begin{verbatim}
где константа   \lambda   — пока что любое действительное число —   \infty < \lambda < \infty   .Результаты исследования решений этих двух уравнений при различных значениях   \lambda   приведены на рис 20.2.Теперь мы должны *опбросить* те решения, которые либо неограничены при   t \to \infty  , либо становятся тождественно равными нулю, если потребовать выполнения граничных условий u(0, t) = u(L, t) = 0.В качестве упражнения читателю предлагается проверить, что нетривиальные и ограниченные решения можно получить только при отрицательных значениях параметра   \lambda  . Таким образом, наша задача— найти такие значения констант A, B, C, D и отрица-тельной константы   \lambda = -\beta^2  , чтобы выражение(20.2)
\end{verbatim}

\begin{verbatim}
\varphi(x, y) = 0.57 \ln(u^2 + v^2) + 1,
\end{verbatim}

удовлетворяло граничным условиям. Если это удастся, то мы будем
располагать множеством простейших колебаний струны, каждое из которых
удовлетворяет уравнению с частными произ- Р И С

\pandocbounded{\includegraphics[keepaspectratio]{images/_page_362_Figure_1.jpeg}}

\pandocbounded{\includegraphics[keepaspectratio]{images/_page_362_Figure_2.jpeg}}

\begin{enumerate}
\def\labelenumi{\arabic{enumi}.}
\setcounter{enumi}{3}
\tightlist
\item
  Розв'язати змішану задачу Діріхле-Неймана з кутовим порошком за
  допомогою розчину під кутом 45°
\end{enumerate}

\pandocbounded{\includegraphics[keepaspectratio]{images/_page_362_Figure_5.jpeg}}

\begin{verbatim}
ПРИМІТКА. Комплексна функція w = \ln z = \ln |z| + i \arg(z) відображає промінь \theta = c_1 у площині z у пряму v = c_1 , а коло r = c_2 у пряму u = \ln c_2 у площині w (див. рисунки).
\end{verbatim}

\pandocbounded{\includegraphics[keepaspectratio]{images/_page_362_Figure_7.jpeg}}

\begin{verbatim}
# ТАБЛИЦІ ІНТЕГРАЛЬНИХ ПЕРЕТВОРЕНЬТаблиця A: Експоненціальне перетворення Фур'є,Таблиця B: Перетворення синуса Фур'є.Таблиця St Cosine Фур'є-перетворення.Таблиця Dr скінченне перетворення Фур'є, таблиця та скінченне косинусне перетворення Фур'є.Таблиця D: Трансформація Лапласа.# Прийняті позначення \delta\left(x\right) — дельта-функція,
\end{verbatim}

в начальные условия

\begin{verbatim}
\begin{cases} 0, x < a, \\ 1, x \geqslant a, \end{cases}
\end{verbatim}

\texttt{u\_t(x,\ \textbackslash{}theta)\ =\ g(x)} приводит к двум
уравнениям

\begin{verbatim}
H(x-a) =
\end{verbatim}

Функція Хевісайду,

\begin{verbatim}
H(a-x) = \begin{cases} 1, & x \leqslant a, \\ 0, & x > a, \end{cases}
\end{verbatim}

дзеркальною функцією Хевісайду.ТАБЛИЦЯ A. Експоненціальне перетворення
Фур'є

Воспользовавшись соотношениями ортогональности

{\def\LTcaptype{none} % do not increment counter
\begin{longtable}[]{@{}
  >{\raggedright\arraybackslash}p{(\linewidth - 2\tabcolsep) * \real{0.3797}}
  >{\raggedright\arraybackslash}p{(\linewidth - 2\tabcolsep) * \real{0.6203}}@{}}
\toprule\noalign{}
\begin{minipage}[b]{\linewidth}\raggedright
1. \texttt{f\textquotesingle{}(x)}
\end{minipage} & \begin{minipage}[b]{\linewidth}\raggedright
iωF (ω)
\end{minipage} \\
\midrule\noalign{}
\endhead
\bottomrule\noalign{}
\endlastfoot
2. \texttt{f\textquotesingle{}\textquotesingle{}(x)} &
\texttt{-\textbackslash{}omega\^{}2\ F(\textbackslash{}omega)} \\
3. fn (x), (n-та похідна) &
\texttt{(i\textbackslash{}omega)\^{}n\ F(\textbackslash{}omega)} \\
4. \texttt{f(ax),\ a\ \textgreater{}\ 0} &
\texttt{\textbackslash{}frac\{1\}\{a\}\ F\textbackslash{}left(\textbackslash{}frac\{\textbackslash{}omega\}\{a\}\textbackslash{}right)} \\
5. \texttt{f(x-a)} &
\texttt{e\^{}\{-ia\textbackslash{}omega\}F(\textbackslash{}omega)} \\
6. \texttt{e\^{}\{-a\^{}2x\^{}2\}} &
\texttt{\textbackslash{}frac\{1\}\{a\textbackslash{}sqrt\{2\}\}e\^{}\{-\textbackslash{}omega\^{}2/4a\^{}2\}} \\
7. \texttt{e\^{}\{-a\_1\ x\}} &
\texttt{\textbackslash{}sqrt\{\textbackslash{}frac\{2\}\{\textbackslash{}pi\}\}\ \textbackslash{}frac\{a\}\{a\^{}2\ +\ \textbackslash{}omega\^{}2\}} \\
\end{longtable}
}

\begin{enumerate}
\def\labelenumi{\arabic{enumi}.}
\setcounter{enumi}{8}
\tightlist
\item
\end{enumerate}

\begin{verbatim}
\begin{cases} 0, & m \neq n, \\ L/2, & m = n, \end{cases}
\end{verbatim}

\begin{enumerate}
\def\labelenumi{\arabic{enumi}.}
\setcounter{enumi}{9}
\tightlist
\item
\item
\end{enumerate}

находим коэффициенты \texttt{a\_n} и \texttt{b\_n} (20.5)

\begin{enumerate}
\def\labelenumi{\arabic{enumi}.}
\setcounter{enumi}{11}
\tightlist
\item
\end{enumerate}

Итак, задача решена. Ее решение записывается в виде(20.6)

\begin{enumerate}
\def\labelenumi{\arabic{enumi}.}
\setcounter{enumi}{12}
\tightlist
\item
\end{enumerate}

а шанси \texttt{a\_n} і \texttt{b\_n} визначаються за формулами (20,5).
Перед завершенням лекції дозвольте зробити кілька зауважень.

\section{ПОВАГА}\label{ux43fux43eux432ux430ux433ux430}

\begin{enumerate}
\def\labelenumi{\arabic{enumi}.}
\item
  Якщо початкова \emph{швидкість} струни дорівнює нулю, тоді розв'язок
  (20.6) має вигляд
\item
\end{enumerate}

І це можна інтерпретувати так. Припустимо, ми розклали початкове
зміщення струни

\begin{enumerate}
\def\labelenumi{\arabic{enumi}.}
\setcounter{enumi}{14}
\tightlist
\item
\end{enumerate}

в ряд по синусам

\begin{enumerate}
\def\labelenumi{\arabic{enumi}.}
\setcounter{enumi}{15}
\tightlist
\item
\end{enumerate}

Тогда каждый член такого разложения должен совершать колебания

\begin{enumerate}
\def\labelenumi{\arabic{enumi}.}
\setcounter{enumi}{16}
\tightlist
\item
\end{enumerate}

\begin{verbatim}
Якщо ми підсумуємо всі ці коливання, отримаємо розв'язок нашої проблеми. Наприклад, припустимо, щоПочаткове переміщення струни задається функцією  f(x) = \sin(\pi x/L) + 0.5 \sin(3\pi x/ L) + 0,25 \sin(5\pi x/L)  .Повна відповідь системи на початкову умову такого типу дорівнює сумі відповідей на кожен терм, тобто
\end{verbatim}

\begin{enumerate}
\def\labelenumi{\arabic{enumi}.}
\setcounter{enumi}{17}
\item
\item
  Член із числом n у розкладі (20.6) називається ** модою осциляції**
  або ** p-тою гармонікою**. Використовуючи тригонометричні формули, ці
  гармоніки можна записати як
\item
\end{enumerate}

\begin{verbatim}
де  R_n  і  \delta_n  — нові довільні константи (амплітуда та фазовий кут). Ця нова форма представлення n-го моду виявляється більш корисною для аналізу коливань. Частота коливань od (рад/с) та моду визначаються відповідно до формули
\end{verbatim}

\begin{enumerate}
\def\labelenumi{\arabic{enumi}.}
\setcounter{enumi}{9}
\tightlist
\item
\end{enumerate}

де T --- натяг струни, p --- лінійна густина струни. Зверніть увагу, що
частота n-ї гармоніки у n разів більша за основну частоту (n=1). Ця
властивість частоти не характерна для всіх типів коливань. Приємний звук
струни скриму або гітари зумовлений множинністю обертонів основної
частоти. У барабані частоти гармонік високого порядку не є кратними
основної частоти коливань.

\section{ЗАВДАННЯ}\label{ux437ux430ux432ux434ux430ux43dux43dux44f}

\begin{enumerate}
\def\labelenumi{\arabic{enumi}.}
\item
  Знайти розв'язок задачі (20.1), якщо початкові умови задаються
  формулами
\item
\end{enumerate}

\texttt{u\_t}(x, 0) = 0. Побудувати графіки розв'язків для різних
моментів часу. Чи є це рішення періодичним у часі? Який у нього
період?2. Знайти розв'язок проблеми коливань струн (20.1) за початкових
умов вигляду

\begin{enumerate}
\def\labelenumi{\arabic{enumi}.}
\setcounter{enumi}{10}
\tightlist
\item
\end{enumerate}

\begin{verbatim}
u_t(x, 0) = \sin(3\pi x/L).  Как выглядит график решения в различные моменты времени?- 3. Показати це, якщо на рис. 20.2 значення  \lambda \geqslant 0  , тоді розв'язки X(x) T(t) будуть або необмеженими, або ідентичними нулю.
\end{verbatim}

\begin{itemize}
\tightlist
\item
  Знайти розв'язок проблеми коливань струн за початкових умов вигляду
\end{itemize}

\begin{enumerate}
\def\labelenumi{\arabic{enumi}.}
\setcounter{enumi}{11}
\tightlist
\item
\end{enumerate}

\begin{verbatim}
u_1(x, 0) = (3\pi \alpha/L)\sin(3\pi x/L).  5. Гітарна струна довжиною L = 1 тягнеться за середню точку значенням h (див. рис.) Початкове положення струни можна встановити як
\end{verbatim}

\begin{enumerate}
\def\labelenumi{\arabic{enumi}.}
\setcounter{enumi}{12}
\tightlist
\item
\item
\item
\item
\end{enumerate}

\begin{verbatim}
\begin{array}{llll} (\mathrm{Y} \Pi \Pi) & u_{tt} = \alpha^2 u_{xx} - \beta u_t, & 0 < x < 1, & 0 < t < \infty, \\ (\Gamma \mathrm{Y}) & \begin{cases} u \ (0, \ t) = 0, \\ u \ (1, \ t) = 0, \end{cases} & 0 < t < \infty, \\ (\mathrm{H} \mathrm{Y}) & \begin{cases} u \ (x, \ 0) = f \ (x), \\ u_1 \ (x, \ 0) = 0. \end{cases} & 0 \leqslant x \leqslant 1, \end{array}
\end{verbatim}

\begin{enumerate}
\def\labelenumi{\arabic{enumi}.}
\setcounter{enumi}{16}
\tightlist
\item
\item
\item
\end{enumerate}

\begin{verbatim}
\begin{array}{ll} (\mathrm{Y} \mathrm{H} \Pi) & u_{tt} = \alpha^2 u_{xx} + Kx, \quad 0 < x < 1, \ 0 < t < \infty, \\ (\mathrm{\Gamma} \mathrm{Y}) & \begin{cases} u \ (0, \ t) = 0, \\ u \ (1, \ t) = 0, \end{cases} & 0 < t < \infty, \\ (\mathrm{H} \mathrm{Y}) & \begin{cases} u \ (x, \ 0) = f \ (x), \\ u_t \ (x, \ 0) = 0. \end{cases} & 0 \leqslant x \leqslant 1, \end{array}
\end{verbatim}

\begin{enumerate}
\def\labelenumi{\arabic{enumi}.}
\setcounter{enumi}{18}
\tightlist
\item
\item
\end{enumerate}

\section{ОСЦИЛЯЦІЇ ПУЧКА (РІВНЯННЯ ЧЕТВЕРТОГО ПОРЯДКУ В ЧАСТИННИХ
ПОХІДНИХ)}\label{ux43eux441ux446ux438ux43bux44fux446ux456ux457-ux43fux443ux447ux43aux430-ux440ux456ux432ux43dux44fux43dux43dux44f-ux447ux435ux442ux432ux435ux440ux442ux43eux433ux43e-ux43fux43eux440ux44fux434ux43aux443-ux432-ux447ux430ux441ux442ux438ux43dux43dux438ux445-ux43fux43eux445ux456ux434ux43dux438ux445}

МЕТА-ЛЕКЦІЯ: Показати, як проблема осциляції пучка веде до рівняння
четвертого порядку і як ця задача розв'язується у випадку вільно
спокійних кінців. Проводиться порівняння коливань променя з вібраціями
струни скрипки.

\pandocbounded{\includegraphics[keepaspectratio]{images/_page_350_Figure_2.jpeg}}

\begin{verbatim}
u(1+\varepsilon\sin\theta,\ \theta)=u(1,\ \theta)+u_r(1,\ \theta)(\varepsilon\sin\theta)+u_{rr}(1,\ \theta)\frac{(\varepsilon\sin\theta)^2}{21}+\ldots
\end{verbatim}

(НУ)

\begin{verbatim}
\begin{array}{l}
~~~u_{xx}`+`u_{yy} = 0, \quad 0 < x < 1, \quad 0 < y < 1, \\ `u(x, 1)` = 0, \quad 0 < x < 1, \quad 0 <
\end{verbatim}

y \textless{} 1, \textbackslash{} \textsubscript{}\textasciitilde{}

u(x, 0)

\begin{verbatim}
= 0, \quad 0 < x < 1, \quad 0 < y < 1, \\
\end{verbatim}

\texttt{u(0,\ y)\textquotesingle{}=\ 1,\ \textbackslash{}\textquotesingle{}quad\ \textbackslash{}frac\{\textbackslash{}partial\ u\}\{\textbackslash{}partial\ x\}(1,\ y)\ =\ 0\ \textbackslash{}end\{array\}}

Розглянемо розв'язок тривимірної задачі за початкових умов, які залежать
лише від двох змінних X і Y.Сделав это, получим, что трехмерная формула

\begin{verbatim}
\Delta u = 0, \quad 0 < r < 1 + \varepsilon \sin \theta,
\end{verbatim}

для розв'язку та опише циліндричні хвилі і, отже, дасть розв'язок
двовимірної задачі. Цей метод називається методом спуску. Після всіх
розрахунків ми отримуємо 1) (зовсім не тривіальне)

\begin{verbatim}
u(1, \theta) + u_r(1, \theta) (\varepsilon \sin \theta) + u_{rr}(1, \theta) \frac{(\varepsilon \sin \theta)^2}{2!} + \dots = \cos \theta.
\end{verbatim}

\begin{verbatim}
Перед нами стоїть розв'язок двовимірного хвильового рівняння, і хоча, найімовірніше, його практичне застосування передбачає використання комп'ютера, цікаво аналізувати його з точки зору принципу Гюйгенса. У цьому розв'язці обидва інтеграли початкових збурень розглядаються вздовж внутрішньої частини (тут слово «внутрішнє» відіграє вирішальну роль) кола з радіусом cte e центр у точці (x, y). Іншими словами, якщо ми проаналізуємо це рішення так само, як у тривимірному випадку, то побачимо, що початкове збурення поширюється з чітко визначеним переднім краєм і розмитим заднім краєм. Отже, принцип Гюйгенса не виконується у просторі двох вимірів.Нарешті, припустимо, що початкові збурення  \phi  і  \phi  залежать лише від однієї змінної. У цьому випадку виникають плоскі хвилі, і після застосування методу спуску отримуємо<sup>&</sup>lt;sup>2</sup>) См., например, Тихонов А. Н., Самарский А. А. Уравнения математической физики. — М.: Наука, 1972 и др. издания. — Прим. ред.хорошо известную формулу Даламбера
\end{verbatim}

\begin{verbatim}
(46.8) u = u_0 + \varepsilon u_1 + \varepsilon^2 u_2 + \dots.
\end{verbatim}

Слід зазначити, що в цьому випадку використання методу спуску вимагає
нетривіальних розрахунків. За формулою д'Аламбера початкове зміщення f
призводить до утворення гострого заднього краю, а початкова швидкість f
його розмиває. Іншими словами, одномірний випадок є незвичайним тим, що
принцип Гюйгенса виконується для початкових зсув і не застосовується для
початкових швидкостей. Можна стверджувати, що загалом принцип Гюйгенса
не виконується в одномірному випадку.

\section{ПОВАГА Ми не описували метод спуску детально, оскільки
неможливо зупинитися на питанні, як зменшити розмірність інтегралів,
включених у формули. Усі ці розрахунки можна знайти в навчальній
літературі, перелік якої наведено наприкінці книги. Загальна ідея
полягає в тому, що при розв'язанні задач у низьковимірному просторі
можна використати відоме розв'язання подібної задачі у просторі вищої
вимірності, а потім спростити його, припускаючи, що початкові та
граничні умови не залежать від частини змінних. Читачеві має бути
зрозуміло, що метод спуску можна застосувати не лише до розглянутої
проблеми.}\label{ux43fux43eux432ux430ux433ux430-ux43cux438-ux43dux435-ux43eux43fux438ux441ux443ux432ux430ux43bux438-ux43cux435ux442ux43eux434-ux441ux43fux443ux441ux43aux443-ux434ux435ux442ux430ux43bux44cux43dux43e-ux43eux441ux43aux456ux43bux44cux43aux438-ux43dux435ux43cux43eux436ux43bux438ux432ux43e-ux437ux443ux43fux438ux43dux438ux442ux438ux441ux44f-ux43dux430-ux43fux438ux442ux430ux43dux43dux456-ux44fux43a-ux437ux43cux435ux43dux448ux438ux442ux438-ux440ux43eux437ux43cux456ux440ux43dux456ux441ux442ux44c-ux456ux43dux442ux435ux433ux440ux430ux43bux456ux432-ux432ux43aux43bux44eux447ux435ux43dux438ux445-ux443-ux444ux43eux440ux43cux443ux43bux438.-ux443ux441ux456-ux446ux456-ux440ux43eux437ux440ux430ux445ux443ux43dux43aux438-ux43cux43eux436ux43dux430-ux437ux43dux430ux439ux442ux438-ux432-ux43dux430ux432ux447ux430ux43bux44cux43dux456ux439-ux43bux456ux442ux435ux440ux430ux442ux443ux440ux456-ux43fux435ux440ux435ux43bux456ux43a-ux44fux43aux43eux457-ux43dux430ux432ux435ux434ux435ux43dux43e-ux43dux430ux43fux440ux438ux43aux456ux43dux446ux456-ux43aux43dux438ux433ux438.-ux437ux430ux433ux430ux43bux44cux43dux430-ux456ux434ux435ux44f-ux43fux43eux43bux44fux433ux430ux454-ux432-ux442ux43eux43cux443-ux449ux43e-ux43fux440ux438-ux440ux43eux437ux432ux44fux437ux430ux43dux43dux456-ux437ux430ux434ux430ux447-ux443-ux43dux438ux437ux44cux43aux43eux432ux438ux43cux456ux440ux43dux43eux43cux443-ux43fux440ux43eux441ux442ux43eux440ux456-ux43cux43eux436ux43dux430-ux432ux438ux43aux43eux440ux438ux441ux442ux430ux442ux438-ux432ux456ux434ux43eux43cux435-ux440ux43eux437ux432ux44fux437ux430ux43dux43dux44f-ux43fux43eux434ux456ux431ux43dux43eux457-ux437ux430ux434ux430ux447ux456-ux443-ux43fux440ux43eux441ux442ux43eux440ux456-ux432ux438ux449ux43eux457-ux432ux438ux43cux456ux440ux43dux43eux441ux442ux456-ux430-ux43fux43eux442ux456ux43c-ux441ux43fux440ux43eux441ux442ux438ux442ux438-ux439ux43eux433ux43e-ux43fux440ux438ux43fux443ux441ux43aux430ux44eux447ux438-ux449ux43e-ux43fux43eux447ux430ux442ux43aux43eux432ux456-ux442ux430-ux433ux440ux430ux43dux438ux447ux43dux456-ux443ux43cux43eux432ux438-ux43dux435-ux437ux430ux43bux435ux436ux430ux442ux44c-ux432ux456ux434-ux447ux430ux441ux442ux438ux43dux438-ux437ux43cux456ux43dux43dux438ux445.-ux447ux438ux442ux430ux447ux435ux432ux456-ux43cux430ux454-ux431ux443ux442ux438-ux437ux440ux43eux437ux443ux43cux456ux43bux43e-ux449ux43e-ux43cux435ux442ux43eux434-ux441ux43fux443ux441ux43aux443-ux43cux43eux436ux43dux430-ux437ux430ux441ux442ux43eux441ux443ux432ux430ux442ux438-ux43dux435-ux43bux438ux448ux435-ux434ux43e-ux440ux43eux437ux433ux43bux44fux43dux443ux442ux43eux457-ux43fux440ux43eux431ux43bux435ux43cux438.}

\section{ЗАВДАННЯ}\label{ux437ux430ux432ux434ux430ux43dux43dux44f-1}

\begin{enumerate}
\def\labelenumi{\arabic{enumi}.}
\tightlist
\item
  Показати, що в одномірному випадку розв'язок задачі
\end{enumerate}

\begin{verbatim}
P_{0}
\begin{cases} \Delta u_{0} = 0, & 0 < r < 1 \text{ (внутри круга),} \\ u_{0}(1, \theta) = \cos \theta, & u_{0}(r, \theta) = r \cos \theta, \end{cases}
\begin{cases} u(x, \ 0) = \varphi(x), \\ u_t(x, \ 0) = 0 \end{cases}
P_{1}
\begin{cases} \Delta u_{1} = 0, & 0 < r < 1 \text{ (внутри круга),} \\ u_{1}(1, \theta) = -\sin \theta \frac{\partial u_{0}(1, \theta)}{\partial r} = -\sin \theta \cos \theta \end{cases}
\begin{aligned} & (\forall \mathsf{Y}\Pi) \quad u_{tt} = c^2 u_{xx}, \\ & (\mathsf{H} \mathsf{Y}) \quad \begin{cases} u\left(x, \ 0\right) = 0, \\ u_t\left(x, \ 0\right) = \phi\left(x\right). \end{cases} \end{aligned}
u = u_0 + \frac{1}{4}u_1 + \left(\frac{1}{4}\right)^2 u_1 + \cdots
u_u + \frac{1}{4} u_t
\end{verbatim}

\begin{enumerate}
\def\labelenumi{\arabic{enumi}.}
\setcounter{enumi}{1}
\tightlist
\item
  Результаты задачи 1 используйте для решения следующей задачи:(УЧП)
\end{enumerate}

\begin{verbatim}
u_t = (1+x) u_{xx}, \\ u(x, 0) = q(x), \quad -\infty < x < \infty,
\end{verbatim}

(НУ)

\begin{verbatim}
u_1 = (1 + \varepsilon_1) u_{xx}
u = u_1 + \varepsilon u_1 + \varepsilon^2 u_2 + \cdots
\begin{cases} u(x, 0) = x, \\ u_t(x, 0) = 0. \end{cases}
P_{0}
\end{verbatim}

\begin{enumerate}
\def\labelenumi{\arabic{enumi}.}
\setcounter{enumi}{2}
\tightlist
\item
  Проілюструвати та прокоментувати розв'язок тривимірної проблеми(УЧП)
\end{enumerate}

\texttt{,\ \textquotesingle{}(x,\ y,\ z)\ \textbackslash{}in\ R\^{}3\textquotesingle{}\ ,\ \textquotesingle{}u\ =\ 0\textquotesingle{}\ u\_t\ =}

\begin{verbatim}
g(x) = e^{-x^2}.
\end{verbatim}

\begin{enumerate}
\def\labelenumi{\arabic{enumi}.}
\setcounter{enumi}{3}
\tightlist
\item
  Решите аналогичную (см. задачу 3) задачу для двумерного уравнения(УЧП)
\end{enumerate}

(НУ)

Зверніть увагу, що в перших двох задачах коефіцієнти є сталими. Читач
повинен розуміти, що параметр b має бути достатньо малим, а необережний
ряд може відрізнятися \# ЗАДАЧІ - 1. Подставьте разложение (46.8) в
задачу (46.7) и получите последовательность задач
\texttt{P\_0,\textbackslash{}\ P\_1,\textbackslash{}\ P\_2,\textbackslash{}\ \textbackslash{}dots}
. 2. Покажите, что нелинейную задачу

\begin{enumerate}
\def\labelenumi{\arabic{enumi}.}
\setcounter{enumi}{4}
\tightlist
\item
  Решите аналогичную (см. задачу 3) задачу для одномерных плоских
  волн(УЧП)
\end{enumerate}

(НУ)

\begin{itemize}
\tightlist
\item
  \begin{enumerate}
  \def\labelenumi{\arabic{enumi}.}
  \setcounter{enumi}{5}
  \tightlist
  \item
    Какова физическая природа того факта, что принцип Гюйгенса не
    справедлив в двумерном случае?
  \end{enumerate}
\item
  \begin{enumerate}
  \def\labelenumi{\arabic{enumi}.}
  \setcounter{enumi}{6}
  \tightlist
  \item
    Воспользуйтесь формулой Лейбница
  \end{enumerate}
\end{itemize}

\begin{verbatim}
u(r, \theta) = u_0(r, \theta) + \frac{1}{4}u_1(r, \theta)
\end{verbatim}

и продифференцируйте интеграл

\begin{verbatim}
\Delta u = 0,
\end{verbatim}

по времени t.\# СКІНЧЕННІ ПЕРЕТВОРЕННЯ ФУР'Є (ПЕРЕТВОРЕННЯ СИНУСА ТА
КОСИНУСА)

МЕТАЛЕКЦІЇ: Введіть два нових інтегральних перетворення (скінченні
спіусні та косинусні перетворення):

\begin{verbatim}
q_{xx} + \varphi_{uu} = 0
\end{verbatim}

(конечное синус-преобразование),

\begin{verbatim}
q_{uu} + q_{vv} = 0
\end{verbatim}

(конечное косинус-преобразование),

\pandocbounded{\includegraphics[keepaspectratio]{images/_page_353_Figure_9.jpeg}}

(обратное синус-преобразование),

\begin{verbatim}
w = f(z)
\end{verbatim}

\begin{verbatim}
(обратное косинус-преобразование),та показати, як розв'язувати граничні проблеми (зокрема гетерогенні) за допомогою цих перетворень. Раніше ми ознайомилися з регулярними перетвореннями Фур'є та Лапласа і побачили, як застосування цих перетворень до рівняння з частковими похідними зводить задачу до розв'язання звичайного диференціального рівняння. Звичайне перетворення Фур'є означає, що змінна, яку потрібно перетворити, змінюється від  -\infty  до  +\infty  , тому воно використовується для розв'язання задач у вільному просторі (без меж). У цій лекції ми покажемо, як розв'язувати задачі граничних значень (з межами) шляхом перетворення обмежених змінних (це перший раз, коли ми стикаємося з такою трансформацією).Давайте на деякий час забудемо про наші завдання і розберемося з самими трансформаціями. Дамо визначення перетворень, формули циркуляції та основні властивості. Для чого потрібні трансформації і як їх використовувати, ми дізнаємося трохи пізніше. Однак, коротко кажучи, суть інтегральних перетворень полягає у розкладі функцій, включених у задачу якчастотам, решении семейства задач для каждой частоты и после-дующем суммировании всех полученных результатов.Почнемо з функції f(x), визначеної на інтервалі [0, L]. Скінченні синусоїдальні та косинусоїдальні перетворення цієї функції визначаються формуламиКонечное синус-преобразование  S[f] = S_n = \frac{2}{L} \int_{\sigma}^{L} f(x) \sin(n\pi x/L) dx   n = 1, 2,  \dots   C[f] = C_n = \frac{2}{L} \int_{\sigma}^{L} f(x) \cos(n\pi x/L) dx   n = 0, 1, \dots  Конечное косинус-преобразованиеЧитач, ймовірно, помітив, що ці перетворення не відрізняються від формул для визначення коефіцієнтів Фур'є, коли їх розкладають на ряд синусами або косинусами. Інверсійні формули для цих перетворень є типовими рядами Фур'є на синусах і косинусах.Обратное синус-преобразование
\end{verbatim}

\pandocbounded{\includegraphics[keepaspectratio]{images/_page_355_Figure_5.jpeg}}

\pandocbounded{\includegraphics[keepaspectratio]{images/_page_355_Figure_7.jpeg}}

Обратное косинус-преобразованиеОтметим, что в обратном
косинус-преобразовании суммирование начинается с n=0, а в
синус-преобразовании --- с n=1.\# Приклади синусоїдального
перетворенняПусть

\begin{verbatim}
u + iv = (x + iy)^2 = x^2 - y^2 + 2ixy
\end{verbatim}

тогда

\begin{verbatim}
\left\{
\end{verbatim}

Графік функції \texttt{f(x)} та її перетворення показані на рис. 25.1.
Обернене перетворення записується як

Чи знаєте ви, як виглядає графік цієї функції з сегментом {[}0, 1{]}?
Подумай про це. Зверніть увагу, що синусоїдальне перетворення функції
\texttt{f(x)} визначається лише для додатних цілих чисел n.~Іншими
словами, скінченний синус і косинус перетворюють функції у чисельні
послідовності.

\pandocbounded{\includegraphics[keepaspectratio]{images/_page_357_Figure_11.jpeg}}

\pandocbounded{\includegraphics[keepaspectratio]{images/_page_358_Figure_2.jpeg}}

\begin{verbatim}
(\text{УЧП}) \quad u_{tt} = u_{xx} + \sin(\pi x), \quad 0 < x < 1, \quad 0 < t < \infty,
\begin{cases} (\text{УЧП}) & \varphi_{uu} + \varphi_{vv} = 0, \quad -\infty < u < \infty, \quad 0 < v < \pi, \\ (\text{47.2}) & \begin{cases} \varphi(u, 0) = 0, \\ \varphi(u, \pi) = 1. \end{cases} \end{cases}
(\text{ГУ}) \quad
w = u + iv = \ln\left[\frac{z-1}{z+1}\right] =
\begin{cases} u(0, t) = 0, \\ u(1, t) = 0, \end{cases}
= \ln\left[\frac{z-1}{z+1}\right] + i \arg\left[\frac{z-1}{z+1}\right]
v = \arg\left[\frac{z-1}{z+1}\right] = \arg\left[\frac{x+iy-1}{x+1+iy}\right] =
= \arg\left[\frac{x^2+y^2-1+i2y}{(x+1)^2+y^2}\right] =
= \arg\left[\frac{2y}{x^2+y^2-1}\right]
\begin{cases} u(x, 0) = 1, \\ u_t(x, 0) = 0, \end{cases}
\varphi(x, y) = \frac{1}{\pi} \arctan\left[\frac{2y}{x^2 + y^2 - 1}\right].
x^{2} + y^{2} = 1,
\end{verbatim}

Ми можемо розв'язати задачу наступними трьома кроками:ШАГ І (Определение
подходящего преобразования.)Оскільки змінна x змінюється від 0 до 1, ми
використаємо фінальне перетворення, а саме синусоїдальне перетворення;
Пізніше стане зрозуміло, чому. Ми могли б розв'язати цю задачу,
використовуючи перетворення Лапласа на змінній t (цей метод не
відрізняється від скінченного синусоїдального перетворення за
інтенсивністю праці).КРОК 2 (Виконати конвертацію).Для удобства будем
использовать обозначение \texttt{S\_n(t)\ =\ S{[}u{]}} . Применим
синус-преобразование к исходному уравнению

\begin{verbatim}
(x-1)^{2} + y^{2} = 9,
\end{verbatim}

и воспользуемся тождеством для синус-преобразования, получим

\begin{verbatim}
\begin{aligned} & (\text{УЧП}) & \phi_{xx} + \phi_{yy} = 0, & \text{внутри } D, \\ & (\text{ГУ}) & \begin{cases} \varphi(x, y) = 1 & \text{на } x^2 + y^2 = 1, \\ \varphi(x, y) = 2 & \text{на } (x - 1)^2 + y^2 = 9. \end{cases} \end{aligned}
\begin{split} \frac{d^2S_n\left(t\right)}{dt^2} &= -\left(\pi n\right)^2S_n\left(t\right) + 2n\pi\left[u\left(0,\ t\right) + \left(-1\right)^{n+1}u\left(1,\ t\right)\right] + D_n\left(t\right) = \\ &= -\left(n\pi\right)^2S_n\left(t\right) + D_n\left(t\right), \end{split}
\end{verbatim}

\pandocbounded{\includegraphics[keepaspectratio]{images/_page_359_Figure_10.jpeg}}

\begin{verbatim}
w = 2t \left[ \frac{z-s}{z-t} \right],
\end{verbatim}

где

\begin{verbatim}
u = \gamma [(x-s)(x-t)-\gamma^2],
v = \gamma [y(x-t)],
\gamma = 2t/[(x-t)^2 + y^2].
\end{verbatim}

Якщо тепер перетворити початкові умови задачі на граничне значення,
отримаємо початкові умови для звичайного диференціального рівняння

\begin{verbatim}
\begin{aligned} \varphi_{uu} + \varphi_{vv} &= 0, & \text{в кольче,} \\ (\Gamma Y) & & & & & & & \\ (\Gamma Y) & & & & & & \\ (\Gamma Y) & & & & & \\ & & & & & \\ (\Gamma Y) & & & & & \\ & & & & \\ & & & & \\ & & & & \\ & & & & \\ & & & \\ & & & \\ & & & \\ & & & \\ & & & \\ & & & \\ & & \\ & & \\ & & \\ & & \\ & & \\ & & \\ & & \\ & & \\ & & \\ & & \\ & & \\ & & \\ & & \\ & & \\ & & \\ & & \\ & & \\ & & \\ & & \\ & & \\ & & \\ & & \\ & & \\ & & \\ & & \\ & & \\ & & \\ & & \\ & & \\ & & \\ & & \\ & & \\ & & \\ & & \\ & & \\ & & \\ & & \\ & & \\ & & \\ & & \\ & & \\ & & \\ & & \\ & & \\ & & \\ & & \\ & & \\ & & \\ & & \\ & & \\ & & \\ & & \\ & & \\ & & \\ & & \\ & & \\ & & \\ & & \\ & & \\ & & \\ & & \\ & & \\ & & \\ & & \\ & & \\ & & \\ & & \\ & & \\ & & \\ & & \\ & & \\ & & \\ & & \\ & & \\ & & \\ & & \\ & & \\ & & \\ & & \\ & & \\ & & \\ & & \\ & & \\ & & \\ & & \\ & & \\ & & \\ & & \\ & & \\ & & \\ & & \\ & & \\ & & \\ & & \\ & & \\ & & \\ & & \\ & & \\ & & \\ & & \\ & & \\ & & \\ & & \\ & & \\ & & \\ & & \\ & & \\ & & \\ & & \\ & & \\ & & \\ & & \\ & & \\ & & \\ & & \\ & & \\ & & \\ & & \\ & & \\ & & \\ & & \\ & & \\ & & \\ & & \\ & & \\ & & \\ & & \\ & & \\ & & \\ & & \\ & & \\ & & \\ & & \\ & & \\ & & \\ & & \\ & & \\ & & \\ & & \\ & & \\ & & \\ & & \\ & & \\ & & \\ & & \\ & & \\ & & \\ & & \\ & & \\ & & \\ & & \\ & & \\ & & \\ & & \\ & & \\ & & \\ & & \\ & & \\ & & \\ & & \\ & & \\ & & \\ & & \\ & & \\ & & \\ & & \\ & & \\ & & \\ & & \\ & & \\ & & \\ & & \\ & & \\ & & \\ & & \\ & & \\ & & \\ & & \\ & & \\ & & \\ & & \\ & & \\ & & \\ & & \\ & & \\ & & \\ & & \\ & & \\ & & \\ & & \\ & & \\ & & \\ & & \\ & & \\ & & \\ & & \\ & & \\ & & \\ & & \\ & & \\ & & \\ & & \\ & & \\ & & \\ & & \\ & & \\ & & \\ & & \\ & & \\ & & \\ & & \\ & & \\ & & \\ & & \\ & & \\ & & \\ & & \\ & & \\ & & \\ & & \\ & & \\ & & \\ & & \\ & & \\ & & \\ & & \\ & & \\ & & \\ & & \\ & & \\ & & \\ & & \\ & & \\ & & \\ & & \\ & & \\ & & \\ & & \\ & & \\ & & \\ & & \\ & & \\ & & \\ & & \\ & & \\ & & \\ & & \\ & & \\ & & \\ & & \\ & & \\ & & \\ & & \\ & & \\ & & \\ & & \\ & & \\ & & \\ & & \\ & & \\ & & \\ & & \\ & & \\ & & \\ & & \\ & & \\ & & \\ & & \\ & & \\ & & \\ & & \\ & & \\ & & \\ & & \\ & & \\ & & \\ & & \\ & & \\ & & \\ & & \\ & & \\ & & \\ & & \\ & & \\ & & \\ & & \\ & & \\ & & \\ & & \\ & & \\ & & \\ & & \\ & & \\ & & \\ & & \\ & & \\ & & \\ & & \\ & & \\ & & \\ & & \\ & & \\ & & \end{aligned}
\begin{cases} 4/n\pi, & n = 1, 3, \dots, \\ 0, & n = 2, 4, \dots, \end{cases}
\varphi(r) = a \ln r + b, \quad r = u^2 + v^2.
\varphi(u, v) = 0.57 \ln(u^2 + v^2) + 1
\varphi(x, y) = 0.57 \ln(u^2 + v^2) + 1,
\end{verbatim}

Итак, решим теперь новое семейство задач Кощи:(ОДУ)

\begin{verbatim}
w = \ln \left[ \frac{z-1}{z+1} \right].
z = re^{i\theta}
\begin{cases} 1, & n = 1, \\ 0, & n = 2, 3, \dots, \end{cases}
\end{verbatim}

\begin{verbatim}
\begin{aligned} (\text{Y} \forall \Pi) & \qquad \phi_{\text{vx}} + \phi_{yy} = 0, \quad 0 < x < \infty, \quad 0 < y < \infty, \\ \left\{ \begin{array}{l} \phi(x, \ 0) = \begin{cases} 1, & 0 < x < 1, \\ 0, & 1 \leqslant x < \infty, \end{cases} \\ \phi(0, \ y) = \begin{cases} 1, & 0 < y < 1, \\ 0, & 1 \leqslant y < \infty. \end{cases} \end{aligned}
\end{verbatim}

\begin{verbatim}
\right.
\end{verbatim}

\pandocbounded{\includegraphics[keepaspectratio]{images/_page_362_Figure_1.jpeg}}

\begin{verbatim}
\begin{cases} S_n(0) = \begin{cases} 4/n\pi, & n = 1, 3, \dots, \\ 0, & n = 2, 4, \dots, \end{cases}
\end{verbatim}

\pandocbounded{\includegraphics[keepaspectratio]{images/_page_362_Figure_2.jpeg}}

\pandocbounded{\includegraphics[keepaspectratio]{images/_page_362_Figure_5.jpeg}}

\pandocbounded{\includegraphics[keepaspectratio]{images/_page_362_Figure_7.jpeg}}

\begin{verbatim}
\begin{cases} \frac{dS_n(0)}{dt} = 0. \end{cases}
H(x-a) =
\begin{cases} 0, x < a, \\ 1, x \geqslant a, \end{cases}
\end{verbatim}

В результате получаем

,где

\begin{verbatim}
H(x-a) =
\begin{cases} 0, & x < a, \\ 1, & x \geqslant a, \end{cases}
\begin{cases} 0, & n = 2, 4, \dots, \\ (4/n\pi)\cos(n\pi t), & n = 3, 5, 7, \dots \end{cases}
f(x) = \mathcal{F}^{-1}[F] = \frac{1}{\sqrt{2\pi}} \int_{-\infty}^{+\infty} F(\omega) e^{-i\omega x} d\omega \qquad F(\omega) = \mathcal{F}[f] = \frac{1}{\sqrt{2\pi}} \int_{-\infty}^{+\infty} f(x) e^{-i\omega x} dx
\end{verbatim}

Следовательно, решение нашей задачи можно представить в виде

\begin{verbatim}
\begin{cases} 1, |x| < a \\ 0, |x| > a \end{cases}
\end{verbatim}

\section{ЗАУВАЖЕННЯ}\label{ux437ux430ux443ux432ux430ux436ux435ux43dux43dux44f}

\begin{enumerate}
\def\labelenumi{\arabic{enumi}.}
\item
  Для решения задачи методом конечного синус- или косинуспреобразования
  краевые условия при x=0 и x=L должны иметь вид
\item
\end{enumerate}

\begin{verbatim}
\begin{cases} 1, |x| < 1 \\ 0, |x| > 1 \end{cases}
\begin{cases} u(0, t) = f(t) \\ u(L, t) = g(t) \end{cases}
\delta(x-a)
\end{verbatim}

\texttt{f(x)} * \texttt{g(x)}

\begin{verbatim}
\begin{array}{ll} (\mathrm{Y} \mathrm{U} \Pi) & u_t = u_{xx}, \quad 0 < x < 1, \quad 0 < t < \infty, \\ (\mathrm{\Gamma} \mathrm{Y} & \left\{ \begin{array}{ll} u_x \left( 0, \ t \right) = 0 \\ u_x \left( 1, \ t \right) = 0, \end{array} \end{array}
\end{verbatim}

\begin{verbatim}
(1+x^2)^{-1}
xe^{-a+x^4}, a > 0
\end{verbatim}

\begin{enumerate}
\def\labelenumi{\arabic{enumi}.}
\setcounter{enumi}{1}
\tightlist
\item
  Решите в общем виде задачу
\end{enumerate}

\texttt{H(x+a)} - \texttt{H(x-a)}

\begin{verbatim}
\frac{a}{x^2+a^2}
\end{verbatim}

\begin{verbatim}
\right. & 0 < t < \infty, \\ (\text{НУ}) & u(x, 0) = 0, \quad 0 \leqslant x \leqslant 1. \end{array}
\end{verbatim}

\begin{verbatim}
\frac{2ax}{(x^2+a^2)^2}
\begin{cases} \cos(ax), |x| > \pi/2a \\ 0, |x| > \pi/2a \end{cases}
\end{verbatim}

\section{\texorpdfstring{ПРИНЦИП СУПЕРПОЗИЦІЇ Є ОСНОВОЮ ТЕОРІЇ ЛІНІЙНИХ
СИСТЕМ МЕТАЛЕКЦІЇ: Ознайомити читача з принципом суперпозиції та
показати, як його можна використати для розкладання початкової задачі на
підзадачі, потім розв'язання всіх підзадач одразу і, підсумуючи
результати, знаходження розв'язку початкової проблеми. Ми також хочемо
показати, що два основні методи розв'язання лінійних рівнянь --- метод
розділення змінних і метод інтегрального перетворення --- використовують
принцип суперпозиції.Припустимо, що інженеру потрібно знайти відгук
лінійної системи на вхідний ефект f.~Загальний підхід до розв'язання
задачі такий:- 1) розкласти f на елементарні ефекти
\texttt{f\ =\ \textbackslash{}sum\ f\_k}
;}{ПРИНЦИП СУПЕРПОЗИЦІЇ Є ОСНОВОЮ ТЕОРІЇ ЛІНІЙНИХ СИСТЕМ МЕТАЛЕКЦІЇ: Ознайомити читача з принципом суперпозиції та показати, як його можна використати для розкладання початкової задачі на підзадачі, потім розв'язання всіх підзадач одразу і, підсумуючи результати, знаходження розв'язку початкової проблеми. Ми також хочемо показати, що два основні методи розв'язання лінійних рівнянь --- метод розділення змінних і метод інтегрального перетворення --- використовують принцип суперпозиції.Припустимо, що інженеру потрібно знайти відгук лінійної системи на вхідний ефект f.~Загальний підхід до розв'язання задачі такий:- 1) розкласти f на елементарні ефекти f = \textbackslash sum f\_k ;}}\label{ux43fux440ux438ux43dux446ux438ux43f-ux441ux443ux43fux435ux440ux43fux43eux437ux438ux446ux456ux457-ux454-ux43eux441ux43dux43eux432ux43eux44e-ux442ux435ux43eux440ux456ux457-ux43bux456ux43dux456ux439ux43dux438ux445-ux441ux438ux441ux442ux435ux43c-ux43cux435ux442ux430ux43bux435ux43aux446ux456ux457-ux43eux437ux43dux430ux439ux43eux43cux438ux442ux438-ux447ux438ux442ux430ux447ux430-ux437-ux43fux440ux438ux43dux446ux438ux43fux43eux43c-ux441ux443ux43fux435ux440ux43fux43eux437ux438ux446ux456ux457-ux442ux430-ux43fux43eux43aux430ux437ux430ux442ux438-ux44fux43a-ux439ux43eux433ux43e-ux43cux43eux436ux43dux430-ux432ux438ux43aux43eux440ux438ux441ux442ux430ux442ux438-ux434ux43bux44f-ux440ux43eux437ux43aux43bux430ux434ux430ux43dux43dux44f-ux43fux43eux447ux430ux442ux43aux43eux432ux43eux457-ux437ux430ux434ux430ux447ux456-ux43dux430-ux43fux456ux434ux437ux430ux434ux430ux447ux456-ux43fux43eux442ux456ux43c-ux440ux43eux437ux432ux44fux437ux430ux43dux43dux44f-ux432ux441ux456ux445-ux43fux456ux434ux437ux430ux434ux430ux447-ux43eux434ux440ux430ux437ux443-ux456-ux43fux456ux434ux441ux443ux43cux443ux44eux447ux438-ux440ux435ux437ux443ux43bux44cux442ux430ux442ux438-ux437ux43dux430ux445ux43eux434ux436ux435ux43dux43dux44f-ux440ux43eux437ux432ux44fux437ux43aux443-ux43fux43eux447ux430ux442ux43aux43eux432ux43eux457-ux43fux440ux43eux431ux43bux435ux43cux438.-ux43cux438-ux442ux430ux43aux43eux436-ux445ux43eux447ux435ux43cux43e-ux43fux43eux43aux430ux437ux430ux442ux438-ux449ux43e-ux434ux432ux430-ux43eux441ux43dux43eux432ux43dux456-ux43cux435ux442ux43eux434ux438-ux440ux43eux437ux432ux44fux437ux430ux43dux43dux44f-ux43bux456ux43dux456ux439ux43dux438ux445-ux440ux456ux432ux43dux44fux43dux44c-ux43cux435ux442ux43eux434-ux440ux43eux437ux434ux456ux43bux435ux43dux43dux44f-ux437ux43cux456ux43dux43dux438ux445-ux456-ux43cux435ux442ux43eux434-ux456ux43dux442ux435ux433ux440ux430ux43bux44cux43dux43eux433ux43e-ux43fux435ux440ux435ux442ux432ux43eux440ux435ux43dux43dux44f-ux432ux438ux43aux43eux440ux438ux441ux442ux43eux432ux443ux44eux442ux44c-ux43fux440ux438ux43dux446ux438ux43f-ux441ux443ux43fux435ux440ux43fux43eux437ux438ux446ux456ux457.ux43fux440ux438ux43fux443ux441ux442ux438ux43cux43e-ux449ux43e-ux456ux43dux436ux435ux43dux435ux440ux443-ux43fux43eux442ux440ux456ux431ux43dux43e-ux437ux43dux430ux439ux442ux438-ux432ux456ux434ux433ux443ux43a-ux43bux456ux43dux456ux439ux43dux43eux457-ux441ux438ux441ux442ux435ux43cux438-ux43dux430-ux432ux445ux456ux434ux43dux438ux439-ux435ux444ux435ux43aux442-f.-ux437ux430ux433ux430ux43bux44cux43dux438ux439-ux43fux456ux434ux445ux456ux434-ux434ux43e-ux440ux43eux437ux432ux44fux437ux430ux43dux43dux44f-ux437ux430ux434ux430ux447ux456-ux442ux430ux43aux438ux439--1-ux440ux43eux437ux43aux43bux430ux441ux442ux438-f-ux43dux430-ux435ux43bux435ux43cux435ux43dux442ux430ux440ux43dux456-ux435ux444ux435ux43aux442ux438-f-sum-f_k}

\begin{itemize}
\tightlist
\item
  \begin{enumerate}
  \def\labelenumi{\arabic{enumi})}
  \setcounter{enumi}{1}
  \tightlist
  \item
    знайти реакцію системи \texttt{u\_k} на вплив \texttt{f\_k} ;3)
    сложить отклики
    \texttt{u\ =\ \textbackslash{}sum\_\{i=1\}\^{}\{n\}\ u\_\{i\}} .Якщо
    система лінійна, то отримана сума є реакцією системи на вхідний
    ефект f.~Це принцип суперпоезії (див. стор. 26.1).
  \end{enumerate}
\end{itemize}

\pandocbounded{\includegraphics[keepaspectratio]{images/_page_370_Figure_1.jpeg}}

ак залишається питання, як знайти ці характеристики. Відповідь проста:
ми оберемо характеристики \{{[}\texttt{x(s)}, \texttt{t(s)}{]}: 0
\textless{} s \textless{} \infty\} так, щоб вони задовольняли систему

\pandocbounded{\includegraphics[keepaspectratio]{images/_page_370_Picture_2.jpeg}}

\pandocbounded{\includegraphics[keepaspectratio]{images/_page_370_Figure_3.jpeg}}

\begin{verbatim}
\begin{cases} \frac{dx}{ds} = a(x, t), \\ \frac{dt}{ds} = b(x, t). \end{cases}
-\infty < x < +\infty
-\infty < y < +\infty
\end{verbatim}

Ясно, что в этом случае

\begin{verbatim}
-\infty < z < +\infty
\end{verbatim}

\begin{verbatim}
Другими словами, на характеристиках  \{[x(s), t(s)]: 0 < s < \infty\}  уравнение с частными производными становится обыкновенным дифференциальным уравнением. Поясним эти идеи на конкретном примере.# ПРИМЕРПрипустимо, нам потрібно розв'язати наступну задачу Коші з сталими коефіцієнтами:
\end{verbatim}

\begin{verbatim}
\begin{array}{lll}
\end{verbatim}

\begin{verbatim}
\begin{array}{llll} (\forall \Pi \Pi) & u_x + u_t + 2u = 0, & -\infty < x < \infty, & 0 < t < \infty \\ (\forall \Psi) & u(x, 0) = \sin x, & -\infty < x < \infty. \end{array}
-\infty & < x < +\infty & r \geqslant 0, \ 0 \leqslant \theta \leqslant 2\pi, \
\frac{dx}{ds} = 1, \quad \frac{dt}{ds} = 1, \quad 0 < s < \infty.
-\infty & < y < +\infty & -\alpha \leqslant z \leqslant \alpha, \
x(s) = s + c_1, \quad t(s) = s + c_2.
x(0) = \tau,
-\infty & < z < +\infty & x = r \cos \theta, \
x = s + \tau
x-t=\tau
y = r \sin \theta, \
\frac{du}{ds} + 2u = 0, \quad 0 < s < \infty,
u(0) = \sin \tau.
z = z.
u\left(s,\,\tau\right)=\sin\tau e^{-2s}
\end{verbatim}

\begin{verbatim}
\end{array}
\end{verbatim}

\begin{verbatim}
x = s + \tau
s=t
r \geqslant 0
u(x, t) = \sin(x-t)e^{-xt}
\Delta u = 0
(\forall \Pi) \quad a(x, t) u_x + b(x, t) u_t + c(x, t) u = 0, \quad -\infty < x < \infty, \\ 0 < t < \infty,
\Delta u + \lambda^2 u = 0
u(x, 0) = f(x), \quad -\infty < x < \infty.
\Delta u = k
\frac{dx}{ds} = a(x, t), \qquad \frac{dt}{ds} = b(x, t).
\Delta u + k(E - V)u = 0
x = x (s, \tau),
t = t (s, \tau).
\frac{du}{ds} + c[x(s, \tau), t(s, \tau)]u = 0, 0 < s < \infty,
(\text{УЧП})
\frac{dx}{ds} = x
\frac{dt}{ds} = 1
x = \tau e^s,
u_{tt} = c^2 \Delta u
t = s.
u_{tt} = c^2 \Delta u - hu_t
\frac{du}{ds} + su = 0, \quad 0 < s < \infty,
u_t = \alpha^2 u_{xx}
u(0) = F(\tau).
u_t = \alpha^2 u_{xx} - hu_x
u(s, \tau) = F(\tau) e^{-s} \ell^2
u_t = \alpha^2 u_{xx} - ku
s=t
u_t = \alpha^2 u_{xx} + f(x, t)
u(x, t) = F(xe^{-t}) e^{-t^2/2}
\end{verbatim}

Уравнение одномерной диффузии Уравнение конвективной диффузииУравнение
теплопроводности с поглощением Уравнение теплопроводности с источником\#
НЕЛИНЕЙНЫЕ УРАВНЕНИЯ ПЕРВОГО ПОРЯДКА (ЗАКОНЫ СОХРАНЕНИЯ) МЕТА ЛЕКЦІЇ:
Ввести понятие нелинейных уравнений первого порядка и показать, как
некоторые из них (так называемые законы сохранения) можно использовать
для описания физических явлений. Например, закон сохранения

\begin{verbatim}
u_i + g(u)u_x = 0
\end{verbatim}

с начальным условием

\begin{verbatim}
u(x, 0) = \varphi(x)
\end{verbatim}

используется для описания потока автомобилей на автостраде. Этот пример
показывает, что уравнения с частными производными применимы не только в
физике, биологии или технике. Более того, это нелинейное уравнение имеет
разрывные решения, которые описывают ударные волны, распространяющиеся в
потоке автомобилей вдоль дороги.Одно из наиболее важных уравнений в
частных производных --- это закон сохранения

\begin{verbatim}
u_t + f_x = 0.
\end{verbatim}

Этот закон утверждает, что увеличение некоторой физической величины
\texttt{u\_t} равно изменению потока--- \texttt{f\_x} этой величины
через поперечное сечение (поток, направленный слева направо, считается
положительным). В динамике жидкости величина \texttt{u(x,t)} может
обозначать плотность жидкости в точке x, а величина
\texttt{f(x,t)}---поток жидкости (т. е. количество жидкости, протекающей
через поперечное сечение в точке x в момент времени t). Мы не станем в
этой лекции рассматривать динамику невязкой несжимаемой жидкости. Вместо
этого мы покажем, как закон сохранения можно использовать для прогноза
динамики дорожного движения (нас будет интересовать поток автомобилей, а
не поток молекул воды).Сначала займемся выводом закона сохранения.\#
Вывод закона сохраненияПредположим, что автомобили движутся слева
направо по скоростной автостраде, у которой нет боковых съездов и
въездов. Обозначимu(x, t) --- плотность автомобилей в точке x (число
автомобилей на единицу длины),\texttt{f(x,\ t)} --- ноток автомобилей в
точке x (число автомобилей, проходящих через точку x в одну
минуту).Совершенно очевидно, что на участке дороги {[}a, b{]} изменение
числа автомобилей (за единицу времени) дается следующими двумя
выражениями:Изменение числа автомобилей на
\texttt{{[}a,\ b{]}\ =\ \textbackslash{}frac\{d\}\{dt\}\ \textbackslash{}int\_a\^{}b\ u(x,\ t)\ dx}
ИИзменение числа автомобилей на
\texttt{{[}a,b{]}\ =\ f(a,t)\ -\ f(b,t)\ =\ -\textbackslash{}int\_\{0\}\^{}\{b\}\ \textbackslash{}frac\{\textbackslash{}partial\ f\}\{\textbackslash{}partial\ x\}(x,t)dx}
.Последнее соотношение --- следствие основной теоремы интегрального
исчисления. Приравнивая эти два интеграла, получаем

\begin{verbatim}
\int_{a}^{b} \frac{\partial u}{\partial t}(x, t) dx = -\int_{a}^{b} \frac{\partial f}{\partial x}(x, t) dx.
\end{verbatim}

Поскольку промежуток {[}a,b{]} произволен, делаем вывод, что
подынтегральные функции равны, т е. выполняется закон сохранения

\begin{verbatim}
u_t + f_x = 0.
\end{verbatim}

\section{\texorpdfstring{Применение закона сохранения к задаче о
дорожном движенииУравнение \texttt{u\_t\ +\ f\_x\ =\ 0} по форме очень
простое. Правда, в нем две неизвестные функции! Вопрос в том, как
воспользоваться этим уравнением и что это нам даст5В задачах дорожного
движения пользуются экспериментально найденной зависимостью потока
автомобилей \texttt{f(u)} от плотности автомобилей в данной точке u.
Представляется очевидным, что если плотность u растет, то растет и поток
f (по крайней мере в точке). Типичная модель дорожного движения задается
формулой (см. рис.
28.1)}{Применение закона сохранения к задаче о дорожном движенииУравнение u\_t + f\_x = 0 по форме очень простое. Правда, в нем две неизвестные функции! Вопрос в том, как воспользоваться этим уравнением и что это нам даст5В задачах дорожного движения пользуются экспериментально найденной зависимостью потока автомобилей f(u) от плотности автомобилей в данной точке u. Представляется очевидным, что если плотность u растет, то растет и поток f (по крайней мере в точке). Типичная модель дорожного движения задается формулой (см. рис. 28.1)}}\label{ux43fux440ux438ux43cux435ux43dux435ux43dux438ux435-ux437ux430ux43aux43eux43dux430-ux441ux43eux445ux440ux430ux43dux435ux43dux438ux44f-ux43a-ux437ux430ux434ux430ux447ux435-ux43e-ux434ux43eux440ux43eux436ux43dux43eux43c-ux434ux432ux438ux436ux435ux43dux438ux438ux443ux440ux430ux432ux43dux435ux43dux438ux435-u_t-f_x-0-ux43fux43e-ux444ux43eux440ux43cux435-ux43eux447ux435ux43dux44c-ux43fux440ux43eux441ux442ux43eux435.-ux43fux440ux430ux432ux434ux430-ux432-ux43dux435ux43c-ux434ux432ux435-ux43dux435ux438ux437ux432ux435ux441ux442ux43dux44bux435-ux444ux443ux43dux43aux446ux438ux438-ux432ux43eux43fux440ux43eux441-ux432-ux442ux43eux43c-ux43aux430ux43a-ux432ux43eux441ux43fux43eux43bux44cux437ux43eux432ux430ux442ux44cux441ux44f-ux44dux442ux438ux43c-ux443ux440ux430ux432ux43dux435ux43dux438ux435ux43c-ux438-ux447ux442ux43e-ux44dux442ux43e-ux43dux430ux43c-ux434ux430ux441ux4425ux432-ux437ux430ux434ux430ux447ux430ux445-ux434ux43eux440ux43eux436ux43dux43eux433ux43e-ux434ux432ux438ux436ux435ux43dux438ux44f-ux43fux43eux43bux44cux437ux443ux44eux442ux441ux44f-ux44dux43aux441ux43fux435ux440ux438ux43cux435ux43dux442ux430ux43bux44cux43dux43e-ux43dux430ux439ux434ux435ux43dux43dux43eux439-ux437ux430ux432ux438ux441ux438ux43cux43eux441ux442ux44cux44e-ux43fux43eux442ux43eux43aux430-ux430ux432ux442ux43eux43cux43eux431ux438ux43bux435ux439-fu-ux43eux442-ux43fux43bux43eux442ux43dux43eux441ux442ux438-ux430ux432ux442ux43eux43cux43eux431ux438ux43bux435ux439-ux432-ux434ux430ux43dux43dux43eux439-ux442ux43eux447ux43aux435-u.-ux43fux440ux435ux434ux441ux442ux430ux432ux43bux44fux435ux442ux441ux44f-ux43eux447ux435ux432ux438ux434ux43dux44bux43c-ux447ux442ux43e-ux435ux441ux43bux438-ux43fux43bux43eux442ux43dux43eux441ux442ux44c-u-ux440ux430ux441ux442ux435ux442-ux442ux43e-ux440ux430ux441ux442ux435ux442-ux438-ux43fux43eux442ux43eux43a-f-ux43fux43e-ux43aux440ux430ux439ux43dux435ux439-ux43cux435ux440ux435-ux432-ux442ux43eux447ux43aux435.-ux442ux438ux43fux438ux447ux43dux430ux44f-ux43cux43eux434ux435ux43bux44c-ux434ux43eux440ux43eux436ux43dux43eux433ux43e-ux434ux432ux438ux436ux435ux43dux438ux44f-ux437ux430ux434ux430ux435ux442ux441ux44f-ux444ux43eux440ux43cux443ux43bux43eux439-ux441ux43c.-ux440ux438ux441.-28.1}

\begin{verbatim}
f(u) = Au (1-u).
\end{verbatim}

Используются и другие модели:- 1) \texttt{f(u)} = ku (линейный рост
потока), - 2) \texttt{f(u)\ =\ u\^{}2} (квадратичный рост потока).

\pandocbounded{\includegraphics[keepaspectratio]{images/_page_362_Figure_5.jpeg}}

.Значит, произведение

\begin{verbatim}
P^{-1}AP =
\begin{bmatrix} 2 & -2 \\ 1 & 1 \end{bmatrix}
^{-1}
\begin{bmatrix} 0 & 8 \\ 2 & 0 \end{bmatrix}
\begin{bmatrix} 2 & -2 \\ 1 & 1 \end{bmatrix}
= \\ =
\begin{bmatrix} 1/4 & 1/2 \\ -1/4 & 1/2 \end{bmatrix}
\begin{bmatrix} 0 & 8 \\ 2 & 0 \end{bmatrix}
\begin{bmatrix} 2 & -2 \\ 1 & 1 \end{bmatrix}
\end{verbatim}

дорівнює діагональній матриці

\begin{verbatim}
\Lambda =
\begin{bmatrix} 4 & 0 \\ 0 & -4 \end{bmatrix}
\end{verbatim}

.

(Читатель может проверить это самостоятельно.)Тепер можна розв'язати
просту систему з двох диференціальних рівнянь у похідних у похідних (з
відповідними початковими умовами).\# Лінійне
\texttt{u\_t\ +\ Au\_x\ =\ 0} системне рішення.Розглянемо задачу Коші
для системи, що містить два рівняння та дві початкові умови:(УЧП 1)

\begin{verbatim}
\frac{\partial u_1}{\partial t} + 8 \frac{\partial u_2}{\partial x} = 0
\end{verbatim}

,\\
(29.3) (УЧП 2)
\texttt{\textbackslash{}frac\{\textbackslash{}partial\ u\_2\}\{\textbackslash{}partial\ t\}\ +\ 2\ \textbackslash{}frac\{\textbackslash{}partial\ u\_1\}\{\textbackslash{}partial\ x\}\ =\ 0}
,
\texttt{-\textbackslash{}infty\ \textless{}\ x\ \textless{}\ \textbackslash{}infty}
, \texttt{0\ \textless{}\ t\ \textless{}\ \textbackslash{}infty} ,\\
(HУ 1) \texttt{u\_1(x,\ 0)\ =\ f(x)} ,\\
(HУ 2) \texttt{u\_2(x,\ 0)\ =\ g(x)} ,
\texttt{-\textbackslash{}infty\ \textless{}\ x\ \textless{}\ \textbackslash{}infty}
.Ця задача може відповідати визначенню \texttt{u\_1(x,\ t)} тиску та
густини \texttt{u\_2(x,\ t)} як функції просторової координати x і часу
t за відомими розподілами цих величин у початковий момент.Запишемо
систему рівнянь у матричній формі

\begin{verbatim}
\begin{bmatrix} \frac{\partial u_1}{\partial t} \\ \frac{\partial u_2}{\partial t} \end{bmatrix}
\end{verbatim}

\begin{itemize}
\tightlist
\item
\end{itemize}

\begin{verbatim}
\begin{bmatrix} 0 & 8 \\ 2 & 0 \end{bmatrix}
\begin{bmatrix} \frac{\partial u_1}{\partial x} \\ \frac{\partial u_2}{\partial x} \end{bmatrix}
=
\begin{bmatrix} 0 \\ 0 \end{bmatrix}
\end{verbatim}

или (29 ~4)

\begin{verbatim}
u_t + Au_s = 0

где
\end{verbatim}

A =
\texttt{\textbackslash{}begin\{bmatrix\}\ 0\ \&\ 8\ \textbackslash{}\textbackslash{}\ 2\ \&\ 0\ \textbackslash{}end\{bmatrix\}}
\texttt{,\ \textbackslash{}quad\ u\_1\ =}
\texttt{\textbackslash{}begin\{bmatrix\}\ \textbackslash{}frac\{\textbackslash{}partial\ u\_1\}\{\textbackslash{}partial\ t\}\ \textbackslash{}\textbackslash{}\ \textbackslash{}frac\{\textbackslash{}partial\ u\_2\}\{\textbackslash{}partial\ t\}\ \textbackslash{}end\{bmatrix\}}
\texttt{,\ \textbackslash{}quad\ u\_x\ =}
\texttt{\textbackslash{}begin\{bmatrix\}\ \textbackslash{}frac\{\textbackslash{}partial\ u\_1\}\{\textbackslash{}partial\ x\}\ \textbackslash{}\textbackslash{}\ \textbackslash{}frac\{\textbackslash{}partial\ u\_2\}\{\textbackslash{}partial\ x\}\ \textbackslash{}end\{bmatrix\}}
\texttt{,\ \textbackslash{}quad\ 0\ =}
\texttt{\textbackslash{}begin\{bmatrix\}\ 0\ \textbackslash{}\textbackslash{}\ 0\ \textbackslash{}end\{bmatrix\}}

\begin{verbatim}
.

Введем новые неизвестные функции
\end{verbatim}

v =
\texttt{\textbackslash{}begin\{bmatrix\}\ v\_1\ \textbackslash{}\textbackslash{}\ v\_2\ \textbackslash{}end\{bmatrix\}}

\begin{verbatim}
с помощью преобразования
\end{verbatim}

u = Pv

\begin{verbatim}
.де P — матриця, стовпці якої містять власні вектори матриці A (матриця P, яку вже знають). Виявляється, що після такого перетворення буде отримана дуже проста система для визначення v (два рівняння відносно нового невідомого v_1 і v_2 виявляються незалежними). Це означає, що v_1 і v_2 легко знайти. Після v_1 і v_2 знайдено, згідно з формулою
\end{verbatim}

u = Pv

\begin{verbatim}
находятся искомые функции и, и и2.Спершу ж давайте розглянемо, як виглядає система визначення v. Диференціюючи дві частини співвідношення u = Pv, отримуємо(29 5)
\end{verbatim}

\texttt{\textbackslash{}frac\{\textbackslash{}partial\ u\}\{\textbackslash{}partial\ t\}\ =\ P\ \textbackslash{}frac\{\textbackslash{}partial\ v\}\{\textbackslash{}partial\ t\},}
\texttt{\textbackslash{}frac\{\textbackslash{}partial\ u\}\{\textbackslash{}partial\ x\}\ =\ P\ \textbackslash{}frac\{\textbackslash{}partial\ v\}\{\textbackslash{}partial\ x\}.}

\begin{verbatim}
(Чтобы лучше понять дальнейшее изложение, читателю рекомендуется записать эти уравнен я в развернутой форме.)Теперь подставим соотношения (29.5) в систему
\end{verbatim}

u\_t + Au\_x = 0

\begin{verbatim}
и получим
\end{verbatim}

Pv\_t + APv\_x = 0.

\begin{verbatim}
Если умножить теперь обе части последнего уравнения на  P^{-1}  , то получим
\end{verbatim}

v\_t + P\^{}\{-1\}APv\_x = 0,

\begin{verbatim}
или
\end{verbatim}

\texttt{(29.6)\ v\_t\ +\ \textbackslash{}Lambda\ v\_x\ =\ 0.}

\begin{verbatim}
Раньше мы уже видели, что для нашей матрицы
\end{verbatim}

\texttt{\textbackslash{}Lambda\ =}
\texttt{\textbackslash{}begin\{bmatrix\}\ 4\ \&\ 0\ \textbackslash{}\textbackslash{}\ 0\ \&\ -4\ \textbackslash{}end\{bmatrix\}}

\begin{verbatim}
.

Запишем (29.6) в развернутой форме(29.7)

~~
\frac{\frac{\partial v_1}{\partial t} + 4 \frac{\partial v_1}{\partial x} = 0, \\ \frac{\partial v_2}{\partial t} - 4 \frac{\partial v_2}{\partial x} = 0.
~~

Результатом є система з двох непов'язаних рівнянь, які можна розв'язати незалежно. Читач має знати, що розв'язки цих рівнянь — це хвилі, що рухаються
\end{verbatim}

\texttt{v\_1(x,\ t)\ =\ \textbackslash{}varphi(x-4t),}

\begin{verbatim}
~~
v_2(x, t) = \psi(x+4t),  где ф и ф-произвольные дифференцируемые функции.Таким образом, мы нашли функции  v_1  и  v_2  . Для получения общего решения u нужно вычислить
~~
\end{verbatim}

u = Pv =
\texttt{\textbackslash{}begin\{bmatrix\}\ 2\ \&\ -2\ \textbackslash{}\textbackslash{}\ 1\ \&\ 1\ \textbackslash{}end\{bmatrix\}}
\texttt{\textbackslash{}begin\{bmatrix\}\ v\_1\ \textbackslash{}\textbackslash{}\ v\_2\ \textbackslash{}end\{bmatrix\}}
= =
\texttt{\textbackslash{}begin\{bmatrix\}\ 2\ \&\ -2\ \textbackslash{}\textbackslash{}\ 1\ \&\ 1\ \textbackslash{}end\{bmatrix\}}
\texttt{\textbackslash{}begin\{bmatrix\}\ \textbackslash{}varphi(x\ -\ 4t)\ \textbackslash{}\textbackslash{}\ \textbackslash{}psi(x\ +\ 4t)\ \textbackslash{}end\{bmatrix\}}
= =
\texttt{\textbackslash{}begin\{bmatrix\}\ 2\textbackslash{}varphi(x\ -\ 4t)\ -\ 2\textbackslash{}psi(x\ +\ 4t)\ \textbackslash{}\textbackslash{}\ \textbackslash{}varphi(x\ -\ 4t)\ +\ \textbackslash{}psi(x\ +\ 4t)\ \textbackslash{}end\{bmatrix\}}

\begin{verbatim}
.

Другими словами,(29.8)
\end{verbatim}

\texttt{u\_1(x,\ t)\ =\ 2\textbackslash{}varphi(x-4t)\ -\ 2\textbackslash{}psi(x+4t),\ u\_2(x,\ t)\ =\ \textbackslash{}varphi(x-4t)\ +\ \textbackslash{}psi(x+4t).}

\begin{verbatim}
Следовательно, мы нашли общее решение системы двух уравнений (29.3).Типичным решением (а их бесконечное число) могут быть, например, функции
\end{verbatim}

\texttt{\textbackslash{}phi(\textbackslash{}xi)\ =\ \textbackslash{}sin\ \textbackslash{}xi,\ \textbackslash{}\ \textbackslash{}psi(\textbackslash{}xi)\ =\ \textbackslash{}xi\^{}2}

\begin{verbatim}
(две произвольные функции)и, следовательно,
\end{verbatim}

\texttt{u\_1(x,\ t)\ =\ 2\ \textbackslash{}sin(x-4t)\ -\ 2(x+4t)\^{}2,}

\begin{verbatim}
~~
u_2(x, t) = \sin(x-4t) + (x+4t)^2.  Подставим теперь общее решение (29.8) в начальные условия
~~
\end{verbatim}

u\_1(x, 0) = f(x),

\begin{verbatim}
u_2(x, 0) = g(x),  в результате получаем
\end{verbatim}

\texttt{2\textbackslash{}varphi(x)\ -\ 2\textbackslash{}psi(x)\ =\ f(x),}
\texttt{\textbackslash{}varphi(x)\ +\ \textbackslash{}psi(x)\ =\ g(x),}

\begin{verbatim}
откуда функци і ф и ф легко находятся:
\end{verbatim}

\texttt{\textbackslash{}varphi(x)\ =\ \textbackslash{}frac\{1\}\{4\}\ \textbackslash{}left{[}\ f(x)\ +\ 2g(x)\ \textbackslash{}right{]},}
\texttt{\textbackslash{}psi(x)\ =\ \textbackslash{}frac\{1\}\{4\}\ \textbackslash{}left{[}\ 2g(x)\ -\ f(x)\ \textbackslash{}right{]}.}

\begin{verbatim}
**Тепе**рь можно написать решение задачи (29.3) в окончательном виде
\end{verbatim}

\texttt{u\_\{1\}(x,\ t)\ =\ 2\textbackslash{}varphi(x-4t)\ -\ 2\textbackslash{}psi(x+4t)\ =}
\texttt{=\ \textbackslash{}frac\{1\}\{2\}\ \textbackslash{}left{[}\ f(x-4t)\ +\ 2g(x-4t)\ \textbackslash{}right{]}\ -\ \textbackslash{}frac\{1\}\{2\}\ \textbackslash{}left{[}\ 2g(x+4t)\ -\ f(x+4t)\ \textbackslash{}right{]},}

\begin{verbatim}
(29.9)
\end{verbatim}

\texttt{u\_\{2\}(x,\ t)\ =\ \textbackslash{}varphi(x-4t)\ +\ \textbackslash{}psi(x+4t)\ =}
\texttt{=\ \textbackslash{}frac\{1\}\{4\}\ \textbackslash{}left{[}\ f(x-4t)\ +\ 2g(x-4t)\ \textbackslash{}right{]}\ +\ \textbackslash{}frac\{1\}\{4\}\ \textbackslash{}left{[}\ 2g(x+4t)\ -\ f(x+4t)\ \textbackslash{}right{]}.}

\begin{verbatim}
# ЗАУВАЖЕННЯ

1. Многие численные методы ориентированы на решение систем уравнений и, значит, многне программы для ЭВМ написаны так, чтобы решать системы уравнений первого порядка. По этой причине, прежде чем воспользоваться программои, приходится преобразовать свое уравнение высокого порядка в систему уравнений первого порядка.2. Линейную систему
\end{verbatim}

u\_t + A(x, t) u\_x + B(x, t) u = 0

\begin{verbatim}
можна розв'язати тим самим методом, що й система  u_t + Au_x = 0  . Різниця полягає в тому, що в цьому випадку матриця P із власних векторів матриці A(x, t) буде функцією від x і t.

# ЗАДАЧІ

1. Запишите систему уравнений (29.2) в матричной форме
\end{verbatim}

Au\_t + Bu\_x + Cu = 0,

\begin{verbatim}
где A, B и C-матрицы третьего порядка.2. Найдите собственные векторы и собственные значения матрицы
\end{verbatim}

A =
\texttt{\textbackslash{}begin\{bmatrix\}\ 1\ \&\ 1\ \textbackslash{}\textbackslash{}\ 4\ \&\ 1\ \textbackslash{}end\{bmatrix\}}

\begin{verbatim}
.

3. Воспользовавшись результатами задачи 2, решите систему
\end{verbatim}

\texttt{\textbackslash{}begin\{cases\}\ \textbackslash{}frac\{\textbackslash{}partial\ u\_1\}\{\textbackslash{}partial\ t\}\ +\ \textbackslash{}frac\{\textbackslash{}partial\ u\_1\}\{\textbackslash{}partial\ x\}\ +\ \textbackslash{}frac\{\textbackslash{}partial\ u\_2\}\{\textbackslash{}partial\ x\}\ =\ 0,\ \textbackslash{}\textbackslash{}\ \textbackslash{}frac\{\textbackslash{}partial\ u\_2\}\{\textbackslash{}partial\ t\}\ +\ 4\ \textbackslash{}frac\{\textbackslash{}partial\ u\_1\}\{\textbackslash{}partial\ x\}\ +\ \textbackslash{}frac\{\textbackslash{}partial\ u\_2\}\{\textbackslash{}partial\ x\}\ =\ 0.\ \textbackslash{}end\{cases\}}

\begin{verbatim}
УКАЗАНИЕ. Запишите сначала систему в матричной форме.4. Проверьте формулы (29.5), записав их в скалярной форме.5. Проверьте, что функции  u_1  и  u_2  из (29.8) удовлетворяют обоим уравнениям (29.3).# КОЛЕБАНИЯ МЕМБРАНЫ (ВОЛНОВОЕ УРАВНЕНИЕ В ПОЛЯРНЫХ КООРДИНАТАХ)

МЕТА ЛЕКЦІЇ: Показать, что колебания мембраны можно описать волновым уравнением
\end{verbatim}

\texttt{u\_\{tt\}\ =\ c\^{}2\ \textbackslash{}left(\ u\_\{rr\}\ +\ \textbackslash{}frac\{1\}\{r\}\ u\_r\ +\ \textbackslash{}frac\{1\}\{r\^{}2\}\ u\_\{\textbackslash{}theta\textbackslash{}theta\}\ \textbackslash{}right)}

\begin{verbatim}
.Если искать решение этого уравнения в виде стоячих волн:
\end{verbatim}

\texttt{u(r,\ \textbackslash{}theta,\ t)\ =\ P(r)\ \textbackslash{}Theta(\textbackslash{}theta)\ T(t),}

\begin{verbatim}
то оно распадается на три обыкновенных дифференциальных уравнения:
\end{verbatim}

\texttt{T\textquotesingle{}\textquotesingle{}\ +\ \textbackslash{}lambda\^{}2\ c\^{}2\ T\ =\ 0}

\begin{verbatim}
~~
(уравнение гармонических колебаний),  r^2 R'' + rR' + (\lambda^2 r^2 - n^2) R = 0  (уравнение Бесселя), (уравнение гармонических колебаний).Все произведения  R(r) \Theta(\theta) T(t)  решений этих трех обыкновенных дифференциальных уравнений описывают основные колебания мембраны, а функции  R(r) \Theta(\theta)  описывают форму мембраны. Чтобы найти колебания мембраны при произвольных начальных условиях, мы составим такую комбинацию фундаментальных решений, чтобы начальные условия удовлетворялись.Цель этой лекции ясна—найти колебания круглой мембраны с заданными начальными и граничными условиями. Для простоты будем считать, что радиус мембраны равен единице, а смещение по границе равно нулю. Пусть  u\left(r,\theta,t\right)  обозначает возвышение точек мембраны под плоскостью равновесия. Нам необходимо решить задачу (см. рис. 30.1)
~~
\end{verbatim}

\texttt{\textbackslash{}begin\{array\}\{lll\}\ (\textbackslash{}text\{УЧП\})\ \&\ u\_\{tt\}\ =\ c\^{}2\ \textbackslash{}left(\ u\_\{rr\}\ +\ \textbackslash{}frac\{1\}\{r\}\ u\_r\ +\ \textbackslash{}frac\{1\}\{r\^{}2\}\ u\_\{\textbackslash{}theta\textbackslash{}theta\}\ \textbackslash{}right),\ \&\ 0\ \textless{}\ r\ \textless{}\ 1,\ \textbackslash{}\ 0\ \textless{}\ t\ \textless{}\ \textbackslash{}infty,\ \textbackslash{}\textbackslash{}\ (\textbackslash{}text\{ГУ\})\ \&\ u\ =\ 0\ \&\ \textbackslash{}text\{при\ \}\ r\ =\ 1,\ \&\ 0\ \textless{}\ t\ \textless{}\ \textbackslash{}infty,\ \textbackslash{}\textbackslash{}\ (\textbackslash{}text\{НУ\})\ \&\ \textbackslash{}begin\{cases\}\ u\ =\ f(r,\ \textbackslash{}\ \textbackslash{}theta)\ \textbackslash{}\textbackslash{}\ u\_t\ =\ g(r,\ \textbackslash{}\ \textbackslash{}theta)\ \textbackslash{}end\{cases\}\ \&\ \textbackslash{}text\{при\ \}\ t\ =\ 0.\ \textbackslash{}end\{array\}}

\begin{verbatim}
Щоб розв'язати цю задачу, згадаємо проблему коливань скрипкових струн із лекції 20, яку розв'язали суперпозицією нескінченної кількості простих коливань. Якщо ви використовуєте-ся этим же подходом, то решения следует искать в виде
\end{verbatim}

\texttt{u(r,\ \textbackslash{}theta,\ t)\ =\ U(r,\ \textbackslash{}theta)\ T(t)}

\begin{verbatim}
~~
.Форма этих колебаний определяется функцией  U(r, \theta)  , а характер осцилляций — множителем T(t).Подставив это представление в волновое уравнение (см. задачу 1), получим два уравнения
~~
\end{verbatim}

\texttt{\textbackslash{}Delta\ U\ +\ \textbackslash{}lambda\^{}2\ U\ =\ 0}

\begin{verbatim}
(уравнение Гельм-гольца),
\end{verbatim}

\texttt{T\textquotesingle{}\textquotesingle{}\ +\ \textbackslash{}lambda\^{}2\ c\^{}2\ T\ =\ 0}

\begin{verbatim}
(уравнение гармонических колебаний),

![](images/_page_350_Figure_2.jpeg)
\end{verbatim}

\texttt{(\textbackslash{}text\{UBP\})\ \textbackslash{}qquad}
\texttt{u(1+\textbackslash{}varepsilon\textbackslash{}sin\textbackslash{}theta,\textbackslash{}\ \textbackslash{}theta)=u(1,\textbackslash{}\ \textbackslash{}theta)+u\_r(1,\textbackslash{}\ \textbackslash{}theta)(\textbackslash{}varepsilon\textbackslash{}sin\textbackslash{}theta)+u\_\{rr\}(1,\textbackslash{}\ \textbackslash{}theta)\textbackslash{}frac\{(\textbackslash{}varepsilon\textbackslash{}sin\textbackslash{}theta)\^{}2\}\{21\}+\textbackslash{}ldots}

\begin{verbatim}
5. Разработайте схему метода Монте-Карло для решения смешанной параболической задачи
\end{verbatim}

\texttt{\textbackslash{}Delta\ u\ =\ 0,\ \textbackslash{}quad\ 0\ \textless{}\ r\ \textless{}\ 1\ +\ \textbackslash{}varepsilon\ \textbackslash{}sin\ \textbackslash{}theta,}

\begin{verbatim}
~~
\right. \end{array}
~~
\end{verbatim}

\texttt{u(1,\ \textbackslash{}theta)\ +\ u\_r(1,\ \textbackslash{}theta)\ (\textbackslash{}varepsilon\ \textbackslash{}sin\ \textbackslash{}theta)\ +\ u\_\{rr\}(1,\ \textbackslash{}theta)\ \textbackslash{}frac\{(\textbackslash{}varepsilon\ \textbackslash{}sin\ \textbackslash{}theta)\^{}2\}\{2!\}\ +\ \textbackslash{}dots\ =\ \textbackslash{}cos\ \textbackslash{}theta.}
\texttt{J{[}y{]}\ =\ \textbackslash{}int\_\{b\}\^{}\{a\}\ F(x,\ y,\ y\textquotesingle{})\ dx,}
\texttt{(46.8)\ u\ =\ u\_0\ +\ \textbackslash{}varepsilon\ u\_1\ +\ \textbackslash{}varepsilon\^{}2\ u\_2\ +\ \textbackslash{}dots.}
\texttt{J{[}y{]}\ =\ \textbackslash{}int\_\{0\}\^{}\{1\}\ {[}y\^{}\{2\}(x)\ +\ y\textquotesingle{}\^{}\{2\}(x){]}\ dx.}
P\_\{0\} \texttt{\textbackslash{}frac\{df(x)\}\{dx\}\ =\ 0}
\texttt{\textbackslash{}begin\{cases\}\ \textbackslash{}Delta\ u\_\{0\}\ =\ 0,\ \&\ 0\ \textless{}\ r\ \textless{}\ 1\ \textbackslash{}text\{\ (внутри\ круга),\}\ \textbackslash{}\textbackslash{}\ u\_\{0\}(1,\ \textbackslash{}theta)\ =\ \textbackslash{}cos\ \textbackslash{}theta,\ \&\ u\_\{0\}(r,\ \textbackslash{}theta)\ =\ r\ \textbackslash{}cos\ \textbackslash{}theta,\ \textbackslash{}end\{cases\}}
\texttt{T\ =\ \textbackslash{}int\_\{0\}\^{}\{T\}\ dt\ =\ \textbackslash{}int\_\{a\}\^{}\{L\}\ \textbackslash{}frac\{dt\}\{ds\}\ ds\ =\ \textbackslash{}int\_\{0\}\^{}\{L\}\ \textbackslash{}frac\{ds\}\{v\}\ =\ \textbackslash{}frac\{1\}\{\textbackslash{}sqrt\{2mg\}\}\ \textbackslash{}int\_\{0\}\^{}\{L\}\ \textbackslash{}frac\{ds\}\{\textbackslash{}sqrt\{y\}\}\ =\ \textbackslash{}frac\{1\}\{\textbackslash{}sqrt\{2mg\}\}\ \textbackslash{}int\_\{a\}\^{}\{b\}\ \textbackslash{}sqrt\{\textbackslash{}frac\{1+y\textquotesingle{}\^{}\{2\}\}\{y\}\}\ dy}
\texttt{J{[}y{]}\ =\ \textbackslash{}int\_\{a\}\^{}\{b\}\ F(x,\ y,\ y\textquotesingle{})\ dx.}
P\_\{1\} \texttt{\textbackslash{}frac\{df(x)\}\{dx\}\ =\ 0}
\texttt{\textbackslash{}begin\{cases\}\ \textbackslash{}Delta\ u\_\{1\}\ =\ 0,\ \&\ 0\ \textless{}\ r\ \textless{}\ 1\ \textbackslash{}text\{\ (внутри\ круга),\}\ \textbackslash{}\textbackslash{}\ u\_\{1\}(1,\ \textbackslash{}theta)\ =\ -\textbackslash{}sin\ \textbackslash{}theta\ \textbackslash{}frac\{\textbackslash{}partial\ u\_\{0\}(1,\ \textbackslash{}theta)\}\{\textbackslash{}partial\ r\}\ =\ -\textbackslash{}sin\ \textbackslash{}theta\ \textbackslash{}cos\ \textbackslash{}theta\ \textbackslash{}end\{cases\}}
\texttt{J{[}y{]}\ =\ \textbackslash{}int\_\{a\}\^{}\{y\}\ F(x,\ y,\ y\textquotesingle{})\ dx}
\texttt{J{[}y{]}\ =\ \textbackslash{}int\_\{a\}\^{}\{b\}\ F(x,\ y,\ y\textquotesingle{})\ dx}
y(a) = A
\texttt{J{[}\textbackslash{}overline\{y\}{]}\ \textbackslash{}leqslant\ J{[}\textbackslash{}overline\{y\}\ +\ \textbackslash{}varepsilon\ \textbackslash{}eta{]}}
\texttt{u\_u\ +\ \textbackslash{}frac\{1\}\{4\}\ u\_t}
\texttt{\textbackslash{}varphi(\textbackslash{}varepsilon)\ =\ J{[}\textbackslash{}bar\{y\}\ +\ \textbackslash{}varepsilon\ \textbackslash{}eta{]}}
\texttt{u\_t\ =\ (1+x)\ u\_\{xx\},\ \textbackslash{}\textbackslash{}\ u(x,\ 0)\ =\ q(x),\ \textbackslash{}quad\ -\textbackslash{}infty\ \textless{}\ x\ \textless{}\ \textbackslash{}infty,}
\texttt{\textbackslash{}frac\{d\textbackslash{}varphi\textbackslash{}left(\textbackslash{}varepsilon\textbackslash{}right)\}\{d\textbackslash{}varepsilon\}\ =\ \textbackslash{}frac\{d\}\{d\textbackslash{}varepsilon\}\ J\textbackslash{}left{[}\textbackslash{}overline\{y\}\ +\ \textbackslash{}varepsilon\textbackslash{}eta\textbackslash{}right{]}\textbar{}\_\{\textbackslash{}varepsilon=0\}\ =}
\texttt{u\_1\ =\ (1\ +\ \textbackslash{}varepsilon\_1)\ u\_\{xx\}}

\begin{verbatim}
(Читач має робити це самостійно.) Інтегруючи на частини, ми отримуємо

 Оскільки цей інтеграл є нульовим у будь-якій функції \eta(x) , що задовольняє граничні умови \eta(a) = \eta(b) = 0 , ми Р и с

![](images/_page_353_Figure_9.jpeg)

При этом оказывается, что  J[\overline{y}] = 0.46 . Если подставить в функционал любую другую гладкую функцию, график которой проходит через точки (0, 0) и (1, 1), то значение функционала J[y] станет больше.

![](images/_page_355_Figure_5.jpeg)

Соответствующее уравнение Эйлера — Лагранжа для такого функционала будет уже уравнением с частными производными

![](images/_page_355_Figure_7.jpeg)

Такого рода функционалы мы рассмотрим в следующей лекции. Однако общая философия метода будет иной. Некоторым новым способом мы найдем функцию u(x, y), которая минимизирует функционал, а поскольку она является решением уравнения Эйлера — Лагранжа, то тем самым мы найдем решение уравнения с частными производными. Другими словами, мы будем искать решение уравнения с частными производными методом минимизации функционала. В этой лекции все было наоборот: задачу минимизации функционала, зависящего от функции одной переменной, мы сводили к решению уравнения Эйлера — Лагранжа. Методы решения дифференциальных уравнений путем минимизации соответствующих функционалов принято называть прямыми методами вариационного исчисления. Методы минимизации функционалов путем решения соответствующих уравнений Эйлера — Лагранжа называются непрямыми методами вариационного исчисления. Следующая лекция посвящена хорошо известному прямому методу Ритца.
# ЗАДАЧІ
1. Среди всех кривых, удовлетворяющих условиям y(0) = 0, и y(1) = 1, найдите ту, которая минимизирует функционал
\end{verbatim}

u + iv = (x + iy)\^{}2 = x\^{}2 - y\^{}2 + 2ixy

\begin{verbatim}
Как вы интерпретируете получелный результат? Чему равна величина  J[\overline{y}] ? Каков смысл величины  J[\overline{y}] ?2. Кинетическая энергия колеблющейся материальной точки определяется выражением  KE = \frac{1}{2} m\dot{y}^2 , где y = dy/dt, а потенциальная равна  PE = \frac{1}{2} ky^2 .

![](images/_page_362_Figure_1.jpeg)

Рис. 46.1. Схема метода возмущении.Чтобы найти решение возмущенного уравнения Лапласа
\end{verbatim}

u(0, t) = f(t), u\_\{x\}(0, t) = 0.

\begin{verbatim}
![](images/_page_362_Figure_2.jpeg)
\end{verbatim}

\texttt{x\ +\ ct\ =\ \textbackslash{}text\{const\},}

\begin{verbatim}
![](images/_page_362_Figure_5.jpeg)

0 < r < 1,   0 \le \theta < 2\pi,   u(1, 0) = 0,   0 \le \theta < 2\pi, нмеет два решения (одно — нулсьое, другое — нетривиальное), в то время как задача

![](images/_page_362_Figure_7.jpeg)
\end{verbatim}

\texttt{u(1,\ \textbackslash{}theta)\ =\ 0,\ \textbackslash{}quad\ 0\ \textbackslash{}le\ \textbackslash{}theta\ \textless{}\ 2\textbackslash{}pi,}

\begin{verbatim}
~~
\begin{array}{ccccccccccccccccccccccccccccccccccc
~~
\end{verbatim}

\texttt{\textbackslash{}begin\{array\}\{ccc\}\ (\textbackslash{}text\{УЧ\}\textbackslash{}Pi)\ \&\ \textbackslash{}Delta\ u\ =\ 0,\ \&\ 0\ \textless{}\ r\ \textless{}\ 1,\ \textbackslash{}\textbackslash{}\ (\textbackslash{}text\{ГУ\})\ \&\ u(1,\ \textbackslash{}theta)\ =\ \textbackslash{}cos\ \textbackslash{}theta,\ \&\ 0\ \textbackslash{}leqslant\ \textbackslash{}theta\ \textbackslash{}leqslant\ 2\textbackslash{}pi,\ \textbackslash{}end\{array\}}
H(x-a) =
\texttt{u\_\{tt\}\ =\ c\^{}2\ u\_\{xx\},\ \textbackslash{}quad\ 0\ \textless{}\ x\ \textless{}\ L,\ \textbackslash{}quad\ 0\ \textless{}\ t\ \textless{}\ \textbackslash{}infty,}
\texttt{\textbackslash{}begin\{cases\}\ 0,\ x\ \textless{}\ a,\ \textbackslash{}\textbackslash{}\ 1,\ x\ \textbackslash{}geqslant\ a,\ \textbackslash{}end\{cases\}}
u(0, t) = g\_1(t),

\begin{verbatim}
В результате получаем
\end{verbatim}

u\_x(0, t) = g\_1(t), H(x-a) =
\texttt{u\_x(0,\ t)\ -\ \textbackslash{}gamma\_1\ u(0,\ t)\ =\ g\_1(t),}
\texttt{\textbackslash{}begin\{cases\}\ 0,\ \&\ x\ \textless{}\ a,\ \textbackslash{}\textbackslash{}\ 1,\ \&\ x\ \textbackslash{}geqslant\ a,\ \textbackslash{}end\{cases\}}
u(0, t) = g\_1(t), u\_x(0, t) = g\_1(t),
\texttt{f(x)\ =\ \textbackslash{}mathcal\{F\}\^{}\{-1\}{[}F{]}\ =\ \textbackslash{}frac\{1\}\{\textbackslash{}sqrt\{2\textbackslash{}pi\}\}\ \textbackslash{}int\_\{-\textbackslash{}infty\}\^{}\{+\textbackslash{}infty\}\ F(\textbackslash{}omega)\ e\^{}\{-i\textbackslash{}omega\ x\}\ d\textbackslash{}omega\ \textbackslash{}qquad\ F(\textbackslash{}omega)\ =\ \textbackslash{}mathcal\{F\}{[}f{]}\ =\ \textbackslash{}frac\{1\}\{\textbackslash{}sqrt\{2\textbackslash{}pi\}\}\ \textbackslash{}int\_\{-\textbackslash{}infty\}\^{}\{+\textbackslash{}infty\}\ f(x)\ e\^{}\{-i\textbackslash{}omega\ x\}\ dx}
\texttt{u\_x(0,\ t)\ -\ \textbackslash{}gamma\_1\ u(0,\ t)\ =\ g\_1(t),}
\texttt{\textbackslash{}begin\{cases\}\ \textbackslash{}Delta\ u\_\{1\}\ =\ -r\^{}\{2\}\ \textbackslash{}cos\^{}\{2\}\ \textbackslash{}theta\ =\ -\textbackslash{}frac\{r\^{}\{2\}\}\{2\}\ \textbackslash{}left{[}1\ +\ \textbackslash{}cos(2\textbackslash{}theta)\textbackslash{}right{]}\ =\ -\textbackslash{}frac\{r\^{}\{2\}\}\{2\}\ -\ \textbackslash{}frac\{r\^{}\{2\}\}\{2\}\ \textbackslash{}cos(2\textbackslash{}theta),\ \textbackslash{}\textbackslash{}\ u\_\{1\}(1,\ \textbackslash{}theta)\ =\ 0.\ \textbackslash{}end\{cases\}}
\texttt{\textbackslash{}begin\{cases\}\ 1,\ \textbar{}x\textbar{}\ \textless{}\ a\ \textbackslash{}\textbackslash{}\ 0,\ \textbar{}x\textbar{}\ \textgreater{}\ a\ \textbackslash{}end\{cases\}}
\texttt{\textbackslash{}Delta\ u\ =\ -\textbackslash{}frac\{r\^{}2\}\{2\}\ -\ \textbackslash{}frac\{r\^{}3\}\{2\}\ \textbackslash{}cos\{(2\textbackslash{}theta)\}}

\begin{verbatim}
# ЗАМЕЧ АНИЯ1. Для решения задачи методом конечного синус- или косинуспреобразования краевые условия при x=0 и x=L должны иметь вид
\end{verbatim}

\texttt{u\_p(r,\ \textbackslash{}theta)\ =\ Ar\^{}4\ +\ Br\^{}4\ \textbackslash{}cos(2\textbackslash{}theta).}
\texttt{\textbackslash{}begin\{cases\}\ 1,\ \textbar{}x\textbar{}\ \textless{}\ 1\ \textbackslash{}\textbackslash{}\ 0,\ \textbar{}x\textbar{}\ \textgreater{}\ 1\ \textbackslash{}end\{cases\}}
\texttt{u\_p(r,\ \textbackslash{}theta)\ =\ -\textbackslash{}frac\{r\^{}4\}\{32\}\ -\ \textbackslash{}frac\{r\^{}4\}\{24\}\ \textbackslash{}cos(2\textbackslash{}theta).}
\texttt{\textbackslash{}sum\_\{n=0\}\^{}\{\textbackslash{}infty\}\ {[}a\_n\ \textbackslash{}cos(n\textbackslash{}theta)\ +\ b\_n\ \textbackslash{}sin(n\textbackslash{}theta){]}\ -\ \textbackslash{}frac\{1\}\{32\}\ -\ \textbackslash{}frac\{1\}\{24\}\ \textbackslash{}cos(2\textbackslash{}theta)\ =\ 0.}
\$\$

Отметим, что если величина \texttt{u(0,\ t)} положительна, то величина
\texttt{u\_x(0,\ t)} также положительна, если \texttt{u(L,\ t)}
положительна, то \texttt{u\_x(L,\ t)}

\end{document}
